\documentclass[11pt, a4paper, oneside]{book}

% ── Engine: XeLaTeX ──────────────────────────────────────────────────────────
\usepackage{fontspec}
\setmainfont{Latin Modern Roman}
\setsansfont{Latin Modern Sans}
% Pin the full system DejaVu (v2.37) by path: TeX-trimmed copies drop
% some math glyphs used in the Lean listings (e.g. the divisibility
% bar U+2223, ⟨ ⟩, ≤).
\setmonofont{DejaVuSansMono.ttf}[
  Path=/usr/share/fonts/truetype/dejavu/,
  BoldFont=DejaVuSansMono-Bold.ttf,
  Scale=0.85]

% ── Core packages ────────────────────────────────────────────────────────────
\usepackage[margin=2.5cm]{geometry}
\usepackage{amsmath, amssymb, amsthm, mathtools}
\usepackage{tikz}
\usetikzlibrary{positioning, arrows.meta, calc, decorations.pathmorphing, cd}
\usepackage{tikz-cd}
\usepackage{minted}
\setminted{
  fontsize=\small,
  breaklines=true,
  bgcolor=codebg,
  frame=none,
  framesep=2mm,
  baselinestretch=1.15,
}
\usepackage{xcolor}
\usepackage{hyperref}
\hypersetup{colorlinks=true, linkcolor=deepblue, urlcolor=midblue, citecolor=accentgreen}
\usepackage{microtype}
\usepackage{enumitem}
\usepackage{booktabs}
\usepackage{tcolorbox}
\tcbuselibrary{theorems, skins, breakable}
\usepackage{fancyhdr}
\usepackage{titlesec}
\usepackage{epigraph}

\usepackage{etoolbox}

% ── Colour palette ───────────────────────────────────────────────────────────
\definecolor{deepblue}{HTML}{1a3a5c}
\definecolor{midblue}{HTML}{2e6da4}
\definecolor{lightblue}{HTML}{d6e8f7}
\definecolor{accentgreen}{HTML}{1e7a4e}
\definecolor{lightgreen}{HTML}{d4edda}
\definecolor{accentpurple}{HTML}{6a3fa0}
\definecolor{lightpurple}{HTML}{ede0f7}
\definecolor{accentorange}{HTML}{b85c00}
\definecolor{lightorange}{HTML}{fde8cc}
\definecolor{accentred}{HTML}{a83232}
\definecolor{lightred}{HTML}{f7d6d6}
\definecolor{codebg}{HTML}{f6f8fa}

% ── Theorem-like environments ────────────────────────────────────────────────
\newtcbtheorem[number within=chapter]{defn}{Definition}{
  enhanced, breakable,
  colback=lightblue, colframe=midblue,
  colupper=black, collower=black,
  fonttitle=\bfseries, coltitle=white,
  separator sign={~---},
}{def}

\newtcbtheorem[number within=chapter]{thm}{Theorem}{
  enhanced, breakable,
  colback=lightgreen, colframe=accentgreen,
  colupper=black, collower=black,
  fonttitle=\bfseries, coltitle=white,
  separator sign={~---},
}{thm}

\newtcbtheorem[number within=chapter]{lem}{Lemma}{
  enhanced, breakable,
  colback=lightgreen!60, colframe=accentgreen!70,
  colupper=black, collower=black,
  fonttitle=\bfseries, coltitle=white,
  separator sign={~---},
}{lem}

\newtcbtheorem[number within=chapter]{prop}{Proposition}{
  enhanced, breakable,
  colback=lightgreen!60, colframe=accentgreen!70,
  colupper=black, collower=black,
  fonttitle=\bfseries, coltitle=white,
  separator sign={~---},
}{prop}

\newtcbtheorem[number within=chapter]{cor}{Corollary}{
  enhanced, breakable,
  colback=lightgreen!40, colframe=accentgreen!50,
  colupper=black, collower=black,
  fonttitle=\bfseries, coltitle=white,
  separator sign={~---},
}{cor}

\newtcolorbox{intuition}[1][]{
  enhanced, breakable,
  colback=lightpurple, colframe=accentpurple,
  colupper=black, fonttitle=\bfseries, coltitle=white,
  title={Intuition},
  #1
}

\newtcolorbox{historical}[1][]{
  enhanced, breakable,
  colback=lightorange, colframe=accentorange,
  colupper=black, fonttitle=\bfseries, coltitle=white,
  title={Historical Note},
  #1
}

\newtcolorbox{warning}[1][]{
  enhanced, breakable,
  colback=lightred, colframe=accentred,
  colupper=black, fonttitle=\bfseries, coltitle=white,
  title={Warning},
  #1
}

\newtcolorbox{exercisebox}[1][]{
  enhanced, breakable,
  colback=lightorange!50, colframe=accentorange!70,
  colupper=black, fonttitle=\bfseries, coltitle=white,
  title={Exercise},
  #1
}

\newtcolorbox{leanbox}[1][]{
  enhanced, breakable,
  colback=codebg, colframe=midblue!60,
  colupper=black, fonttitle=\bfseries\ttfamily, coltitle=white,
  title={Lean 4},
  #1
}

% ── Custom macros ────────────────────────────────────────────────────────────
\newcommand{\Z}{\mathbb{Z}}
\newcommand{\Q}{\mathbb{Q}}
\newcommand{\R}{\mathbb{R}}
\newcommand{\C}{\mathbb{C}}
\newcommand{\N}{\mathbb{N}}
\newcommand{\F}{\mathbb{F}}
\newcommand{\GF}{\mathrm{GF}}
\newcommand{\Gal}{\operatorname{Gal}}
\newcommand{\Aut}{\operatorname{Aut}}
\newcommand{\Fix}{\operatorname{Fix}}
\newcommand{\Inv}{\operatorname{Inv}}
% Standard-terminology maps, matching the Lean names fixingGroup /
% fixedField (field->group and group->field respectively).
\newcommand{\fixgrp}{\operatorname{fixingGroup}}
\newcommand{\fixfld}{\operatorname{fixedField}}
\newcommand{\Stab}{\operatorname{Stab}}
\newcommand{\Orb}{\operatorname{Orb}}
\newcommand{\ord}{\operatorname{ord}}
\newcommand{\lcm}{\operatorname{lcm}}
\newcommand{\id}{\operatorname{id}}
\newcommand{\im}{\operatorname{im}}
\renewcommand{\ker}{\operatorname{ker}}
\newcommand{\Hom}{\operatorname{Hom}}
\newcommand{\End}{\operatorname{End}}
\newcommand{\charr}{\operatorname{char}}
\renewcommand{\deg}{\operatorname{deg}}
\newcommand{\Sym}{\mathfrak{S}}
\newcommand{\Alt}{\mathfrak{A}}
\newcommand{\normal}{\trianglelefteq}
\newcommand{\iso}{\cong}
\newcommand{\acts}{\curvearrowright}
\newcommand{\lc}[1]{\texttt{#1}}
\newcommand{\lean}[1]{\mintinline{lean}{#1}}

% Fix: ensure section/chapter title colors don't bleed into body text
\titleformat{\chapter}[display]
  {\normalfont\huge\bfseries\color{deepblue}}
  {\chaptertitlename\ \thechapter}{20pt}{\Huge\color{deepblue}}[\color{black}]
\titleformat{\section}
  {\normalfont\Large\bfseries\color{deepblue}}
  {\thesection}{1em}{}[\color{black}]
\titleformat{\subsection}
  {\normalfont\large\bfseries\color{midblue}}
  {\thesubsection}{1em}{}[\color{black}]

% Global: ensure text is black after any tcolorbox
\AfterEndEnvironment{defn}{\color{black}}
\AfterEndEnvironment{thm}{\color{black}}
\AfterEndEnvironment{lem}{\color{black}}
\AfterEndEnvironment{prop}{\color{black}}
\AfterEndEnvironment{cor}{\color{black}}
\AfterEndEnvironment{intuition}{\color{black}}
\AfterEndEnvironment{historical}{\color{black}}
\AfterEndEnvironment{warning}{\color{black}}
\AfterEndEnvironment{exercisebox}{\color{black}}
\AfterEndEnvironment{leanbox}{\color{black}}
\AfterEndEnvironment{minted}{\color{black}}

% ── Headers/footers ──────────────────────────────────────────────────────────
\pagestyle{fancy}
\fancyhf{}
\fancyhead[L]{\small\itshape\leftmark}
\fancyhead[R]{\small\thepage}
\renewcommand{\headrulewidth}{0.4pt}

% ── Counter for exercises ────────────────────────────────────────────────────
\newcounter{exercise}[chapter]
\renewcommand{\theexercise}{\thechapter.\arabic{exercise}}

% ══════════════════════════════════════════════════════════════════════════════
\begin{document}

% ── Title page ───────────────────────────────────────────────────────────────
\begin{titlepage}
\centering
\vspace*{3cm}
{\Huge\bfseries\color{deepblue} Galois Theory\\[0.3cm] for Programmers\par}
\vspace{1cm}
{\Large\color{midblue} From Symmetries to Unsolvability\\[0.2cm]
in Lean\,4\par}
\vspace{2cm}
{\large A self-contained journey through abstract algebra,\\
from group axioms to the insolvability of the quintic,\\
with every definition formalized and every theorem proved.\par}
\vspace{3cm}
{\large\itshape Draft --- \today\par}
\vfill
\end{titlepage}

% ── Front matter ─────────────────────────────────────────────────────────────
\frontmatter
\tableofcontents

\chapter*{Preface}
\addcontentsline{toc}{chapter}{Preface}

This book tells the most dramatic story in the history of mathematics, and asks
you to formalize it in a proof assistant.

The story begins in the 1770s, when Joseph-Louis Lagrange noticed that the
roots of a polynomial equation have \emph{symmetries}---you can rearrange them
in certain ways without breaking the algebraic relationships between them. Over
the next sixty years, a succession of mathematicians---Gauss, Abel, Galois,
Cauchy---gradually uncovered a deep connection between these symmetries and the
solvability of the equation. The climax came in 1832, when Évariste Galois,
a twenty-year-old political radical in Paris, wrote down the definitive answer
the night before he was killed in a duel: a polynomial equation can be solved
by radicals if and only if its \emph{group of symmetries} has a certain
structural property called \emph{solvability}.

This is not just a historical curiosity. The mathematical framework Galois
created---now called Galois theory---turns out to be fundamental to modern
cryptography (the AES encryption algorithm does arithmetic in a Galois field),
coding theory (Reed-Solomon codes work over finite field extensions), and the
deepest parts of number theory and algebraic geometry.

\paragraph{What this book is.}
A self-contained, zero-to-hero treatment of Galois theory for programmers, with
all definitions, theorems, and proofs formalized in Lean\,4. No prior knowledge
of abstract algebra is assumed. If you can write a function and follow a logical
argument, you have everything you need. Every concept is accompanied by:
\begin{itemize}[nosep]
\item A \emph{historical narrative} explaining where the idea came from and who
      discovered it.
\item A \emph{formal definition} in standard mathematical notation.
\item A \emph{Lean\,4 formalization}: a typeclass, structure, or theorem
      statement that you can type-check.
\item \emph{Worked examples} computed by hand and verified in Lean.
\item \emph{Exercises} ranging from ``prove this basic property'' to
      ``formalize this entire section.''
\end{itemize}

\paragraph{What this book is not.}
A graduate algebra textbook. We sacrifice generality for clarity: we work with
finite groups and finite-degree field extensions, we avoid category theory
(though we point out connections), and we prefer explicit computation over
abstract elegance when the two conflict. The goal is understanding, not
maximal generality.

\paragraph{Prerequisites.}
A little comfort with Lean\,4 syntax helps, but is not assumed: every keyword,
tactic, and construct this book uses---\lc{inductive}, \lc{def}, \lc{structure},
\lc{class}/\lc{instance}, \lc{theorem}, \lc{intro}, \lc{rfl}, \lc{simp},
\lc{rw}, \lc{decide}, \lc{native\_decide}, \lc{omega}, \lc{cases},
\lc{induction}, \lc{obtain}, \lc{calc}, \lc{termination\_by}, and the rest---is
explained from scratch in Appendix~\ref{app:primer}, \emph{A Lean 4 Primer},
each with a small machine-checked example. Dip into it whenever a new construct
appears. No prior abstract algebra.

\paragraph{Lean\,4 setup.}
All code in this book is written for \textbf{vanilla Lean\,4, core library only,
no Mathlib dependency}, and was type-checked against toolchain
\texttt{leanprover/lean4:v4.31.0}. To reproduce it you need nothing more than a
\texttt{lean-toolchain} file containing that one line---no \texttt{lakefile}, no
external packages. We build every algebraic structure from scratch so you see
exactly what is happening inside the typeclasses; where core Lean lacks a
convenience that Mathlib would supply (for instance a \lc{Set} type), we define
our own in a few lines rather than reaching for a library. Every listing in the
book has a compiling counterpart in the companion file
\texttt{GaloisVerification.lean}.

\mainmatter
\color{black}  % Ensure body text starts black

% ══════════════════════════════════════════════════════════════════════════════
\part{The Algebraic Toolkit}
% ══════════════════════════════════════════════════════════════════════════════

% ══════════════════════════════════════════════════════════════════════════════
\chapter{Groups --- The Algebra of Symmetry}
% ══════════════════════════════════════════════════════════════════════════════

\epigraph{``The art of doing mathematics consists in finding that special case
which contains all the germs of generality.''}{\textit{David Hilbert}}

\section{Where Groups Come From}

In the 1770s, Joseph-Louis Lagrange was studying a question that had occupied
mathematicians for centuries: \emph{how do you solve polynomial equations?}

The quadratic formula was ancient knowledge. The cubic and quartic formulas had
been discovered in the 1500s by the Italian mathematicians del Ferro,
Tartaglia, Cardano, and Ferrari---amid a web of secrecy, betrayal, and public
mathematical duels that reads more like a Renaissance thriller than a
mathematics paper. But for degree five and higher, nobody could find a formula.

\begin{historical}[title={The Italian Algebraists}]
Scipione del Ferro secretly discovered the cubic formula around 1515 but told
only his student Antonio Fior. When Niccolò Tartaglia independently found the
formula in 1535, Fior challenged him to a public contest. Tartaglia won
decisively, then shared his secret with Gerolamo Cardano under a sworn oath of
secrecy. Cardano later published it anyway (in his 1545 \textit{Ars Magna}),
claiming del Ferro's priority justified breaking the oath. Tartaglia was
furious. The quartic formula was found by Cardano's student Ludovico Ferrari,
who then challenged Tartaglia to another public contest---and won.

This entire soap opera happened because a general formula for solving equations
was considered the ultimate mathematical achievement. The quest for the quintic
formula would consume mathematicians for the next 250 years.
\end{historical}

Lagrange's approach was novel: instead of trying to find the quintic formula
directly, he asked \emph{why} the existing formulas worked. He noticed that
the key step in solving a cubic or quartic is finding an expression in the
roots that takes on fewer values when the roots are \emph{permuted}. For
example, if a cubic has roots $r_1, r_2, r_3$, the expression
$r_1 + \omega r_2 + \omega^2 r_3$ (where $\omega = e^{2\pi i/3}$) takes
only three values under all six permutations of the roots. This ``reduction
in the number of values'' is what makes the equation solvable.

Lagrange didn't have the word ``group''---that term wouldn't be coined until
Galois used it in 1832---but he was seeing \emph{symmetry}: the permutations
of the roots, and the structure of how those permutations compose with each
other, control whether the equation can be solved.

\begin{intuition}
Think of a group as a mathematical way to talk about \emph{symmetry}. The
symmetries of any object---a geometric shape, a set of roots, a crystal
structure---form a group. The group captures \emph{which rearrangements are
possible} and \emph{how they compose}.
\end{intuition}

\section{The Group Axioms}

After Lagrange, it took nearly a century for the abstract concept of a
``group'' to crystallize. Cauchy (1815) studied permutation groups
systematically. Cayley (1854) gave the first abstract definition. But it was
only in the late 1800s, through the work of Kronecker, Weber, and Frobenius,
that the modern axioms became standard.

\begin{defn}{Group}{group}
A \textbf{group} is a set $G$ together with a binary operation
$\cdot : G \times G \to G$ satisfying:
\begin{enumerate}[nosep, label=(\roman*)]
\item \textbf{Associativity:}
  $(a \cdot b) \cdot c = a \cdot (b \cdot c)$ for all $a, b, c \in G$.
\item \textbf{Identity:} There exists an element $e \in G$ such that
  $e \cdot a = a \cdot e = a$ for all $a \in G$.
\item \textbf{Inverses:} For each $a \in G$, there exists $a^{-1} \in G$
  such that $a \cdot a^{-1} = a^{-1} \cdot a = e$.
\end{enumerate}
If additionally $a \cdot b = b \cdot a$ for all $a, b \in G$, the group is
called \textbf{abelian} (or \textbf{commutative}).
\end{defn}

\begin{historical}[title={Why ``abelian''?}]
Named after Niels Henrik Abel (1802--1829), the Norwegian mathematician who
proved the general quintic is unsolvable by radicals. We will meet Abel's
theorem in Chapter~10. He died of tuberculosis at 26, in poverty, two days
before a letter arrived offering him a prestigious professorship in Berlin.
\end{historical}

\subsection{Lean\,4 Formalization}

We define groups as a typeclass, just as Lean's standard library defines
\lc{Add} or \lc{Mul}. (If any Lean keyword or tactic below is unfamiliar,
Appendix~\ref{app:primer} explains each one with a worked example.)

\begin{leanbox}
\begin{minted}{lean}
class Group (G : Type) where
  mul : G → G → G
  one : G
  inv : G → G
  mul_assoc : ∀ a b c : G, mul (mul a b) c = mul a (mul b c)
  one_mul   : ∀ a : G, mul one a = a
  mul_one   : ∀ a : G, mul a one = a
  inv_mul   : ∀ a : G, mul (inv a) a = one
  mul_inv   : ∀ a : G, mul a (inv a) = one

-- Notation
instance [Group G] : Mul G where mul := Group.mul
instance [Group G] : One G where one := Group.one
instance [Group G] : Inv G where inv := Group.inv
\end{minted}
\end{leanbox}

\begin{intuition}
Notice the parallel with programming interfaces. A \lc{Group G} instance
is exactly an \emph{implementation} of the group interface for the type
\lc{G}. The axioms (\lc{mul\_assoc}, \lc{one\_mul}, etc.) are
\emph{contracts} that any implementation must satisfy---and in Lean, satisfying
a contract means providing a \emph{proof}.
\end{intuition}

\section{First Examples}

\subsection{The Integers Under Addition}

The simplest infinite group: $(\Z, +)$ with identity $0$ and inverse $-n$.

\begin{leanbox}
\begin{minted}{lean}
instance : Group Int where
  mul := (· + ·)
  one := 0
  inv := (- ·)
  mul_assoc := Int.add_assoc
  one_mul   := Int.zero_add
  mul_one   := Int.add_zero
  inv_mul   := Int.add_left_neg
  mul_inv   := Int.add_right_neg
\end{minted}
\end{leanbox}

Every axiom is discharged by a single lemma from Lean's standard library. This
is because the integers were \emph{designed} with these properties in mind---the
group axioms are capturing something that was always there.

\subsection{Modular Arithmetic --- $\Z/n\Z$}

\begin{historical}[title={Gauss and the \textit{Disquisitiones}}]
Carl Friedrich Gauss published his \textit{Disquisitiones Arithmeticae} in 1801,
at the age of 24. It systematized the theory of modular arithmetic (``clock
arithmetic'') and introduced the notation $a \equiv b \pmod{n}$. The
\textit{Disquisitiones} is one of the most influential mathematics books ever
written---it transformed number theory from a collection of scattered results
into a coherent discipline. Gauss didn't use the word ``group,'' but the
structures he studied---integers modulo $n$, primitive roots, quadratic
residues---are all group-theoretic.
\end{historical}

The group $\Z/n\Z$ consists of the integers $\{0, 1, \ldots, n-1\}$ with
addition modulo $n$. In Lean, this is naturally \lc{Fin n}:

\begin{leanbox}
\begin{minted}{lean}
-- We represent ℤ/nℤ as Fin n (for n ≥ 1)
-- Fin n = { val : Nat // val < n }
-- Addition is modular: (a + b) mod n

def Fin.addMod (a b : Fin n) : Fin n :=
  -- `Nat.mod_lt` needs 0 < n; we recover it from `a : Fin n`, whose
  -- proof `a.isLt : a.val < n` forces n to be positive.
  ⟨(a.val + b.val) % n,
    Nat.mod_lt _ (Nat.lt_of_le_of_lt (Nat.zero_le a.val) a.isLt)⟩

def Fin.negMod (a : Fin n) (h : n > 0) : Fin n :=
  ⟨(n - a.val) % n, Nat.mod_lt _ h⟩
\end{minted}
\end{leanbox}

\begin{exercisebox}[title={Exercise 1.1}]
\stepcounter{exercise}
Prove that \lc{Fin n} with \lc{addMod} satisfies the group axioms for any
$n \geq 1$. The hardest part is associativity---you will need properties of
the modular arithmetic operations. \emph{Hint:} use \lc{Nat.add\_mod} and
\lc{omega}.
\end{exercisebox}

\subsection{Symmetries of a Triangle --- $S_3$}

The \textbf{symmetric group} $S_n$ is the group of all permutations of
$n$ elements. $S_3$---the symmetries of $\{1, 2, 3\}$---is the smallest
non-abelian group (6 elements, and the order of multiplication matters).

Think of it as the symmetries of an equilateral triangle: three rotations
(by $0°$, $120°$, $240°$) and three reflections.

\begin{center}
\begin{tikzpicture}[scale=1.2]
  \coordinate (A) at (90:1.5);
  \coordinate (B) at (210:1.5);
  \coordinate (C) at (330:1.5);
  \draw[thick] (A) -- (B) -- (C) -- cycle;
  \node[above] at (A) {$1$};
  \node[below left] at (B) {$2$};
  \node[below right] at (C) {$3$};
  \node[below=1.5cm] {$e$ (identity)};
\end{tikzpicture}
\hspace{2cm}
\begin{tikzpicture}[scale=1.2]
  \coordinate (A) at (90:1.5);
  \coordinate (B) at (210:1.5);
  \coordinate (C) at (330:1.5);
  \draw[thick] (A) -- (B) -- (C) -- cycle;
  \node[above] at (A) {$2$};
  \node[below left] at (B) {$3$};
  \node[below right] at (C) {$1$};
  \draw[->, thick, accentpurple] (0,0) ++(20:0.4) arc (20:100:0.4);
  \node[below=1.5cm] {$r$ (rotate $120°$)};
\end{tikzpicture}
\end{center}

\begin{leanbox}
\begin{minted}{lean}
-- S₃ as an explicit enumeration
inductive S3 where
  | e                    -- identity
  | r                    -- rotation 120°
  | r2                   -- rotation 240°
  | s                    -- reflection (fixes vertex 1)
  | sr                   -- reflection (fixes vertex 3)
  | sr2                  -- reflection (fixes vertex 2)
  deriving DecidableEq, Repr

namespace S3

-- Multiplication table (composition of symmetries: apply the right
-- factor first). This table is associative -- Exercise 1.2 verifies it.
def mul : S3 → S3 → S3
  | e,   x   => x
  | x,   e   => x
  | r,   r   => r2
  | r,   r2  => e
  | r2,  r   => e
  | r2,  r2  => r
  | s,   s   => e
  | s,   r   => sr2
  | s,   r2  => sr
  | r,   s   => sr
  | r2,  s   => sr2
  | sr,  sr  => e
  | sr2, sr2 => e
  | s,   sr  => r2
  | s,   sr2 => r
  | sr,  s   => r
  | sr2, s   => r2
  | r,   sr  => sr2
  | r,   sr2 => s
  | r2,  sr  => s
  | r2,  sr2 => sr
  | sr,  r   => s
  | sr,  r2  => sr2
  | sr2, r   => sr
  | sr2, r2  => s
  | sr,  sr2 => r2
  | sr2, sr  => r

def inv : S3 → S3
  | e   => e
  | r   => r2
  | r2  => r
  | s   => s
  | sr  => sr
  | sr2 => sr2

end S3
\end{minted}
\end{leanbox}

The key observation: $S_3$ is \textbf{not abelian}. We have $r \cdot s = sr$
but $s \cdot r = sr2$, so $r \cdot s \neq s \cdot r$. The order of
operations matters---just as it does for geometric transformations (rotating
then reflecting is different from reflecting then rotating).

\begin{exercisebox}[title={Exercise 1.2}]
\stepcounter{exercise}
Prove that \lc{S3} with the multiplication table above satisfies the group
axioms. Since \lc{S3} is finite with 6 elements, associativity can be proved
by exhaustive case analysis: \lc{decide} or pattern matching on all $6^3 = 216$
triples. \emph{Warning:} this is tedious but instructive. Consider using
\lc{native\_decide} for efficiency.
\end{exercisebox}

\begin{exercisebox}[title={Exercise 1.3}]
\stepcounter{exercise}
Prove that $S_3$ is not abelian by exhibiting a concrete counterexample:
find $a, b \in S_3$ such that $a \cdot b \neq b \cdot a$, and prove the
inequality in Lean.
\end{exercisebox}

\section{Euler and Modular Exponentiation}

\begin{historical}[title={Euler's Theorem}]
Leonhard Euler (1707--1783) is the most prolific mathematician in history,
with over 800 published papers. Among his countless contributions, his work
on modular arithmetic laid the foundations for what we now recognize as group
theory. Euler's theorem---$a^{\varphi(n)} \equiv 1 \pmod{n}$ for $\gcd(a,n)=1$%
---is, in modern language, a consequence of Lagrange's theorem applied to the
multiplicative group $(\Z/n\Z)^\times$. Euler proved it by direct combinatorial
reasoning in 1736, more than a century before the group-theoretic framework
existed to express it cleanly.

Euler's totient function $\varphi(n)$ counts the integers from 1 to $n$ that
are coprime to $n$. In group-theoretic terms, $\varphi(n) = |(\Z/n\Z)^\times|$%
---the order of the multiplicative group modulo $n$.
\end{historical}

We won't prove Euler's theorem yet---it will fall out as an immediate
corollary of Lagrange's theorem in Chapter~2. For now, let's just implement
modular exponentiation, which is the computational heart of RSA encryption:

\begin{leanbox}
\begin{minted}{lean}
-- Fast modular exponentiation by repeated squaring
def modPow (base exp modulus : Nat) : Nat :=
  if modulus ≤ 1 then 0
  else go base exp modulus 1
where
  go (base exp modulus acc : Nat) : Nat :=
    if exp = 0 then acc % modulus
    else
      let base' := (base * base) % modulus
      let acc'  := if exp % 2 = 1 then (acc * base) % modulus else acc
      go base' (exp / 2) modulus acc'
  termination_by exp
\end{minted}
\end{leanbox}

\section{Homomorphisms --- Structure-Preserving Maps}

A \emph{homomorphism} between groups is a function that ``respects'' the group
operation. This is the fundamental notion of ``sameness'' in algebra.

\begin{defn}{Group Homomorphism}{homomorphism}
Let $(G, \cdot_G)$ and $(H, \cdot_H)$ be groups. A function
$\varphi : G \to H$ is a \textbf{homomorphism} if
\[
  \varphi(a \cdot_G b) = \varphi(a) \cdot_H \varphi(b)
  \qquad \text{for all } a, b \in G.
\]
If $\varphi$ is additionally a bijection, it is an \textbf{isomorphism}, and
we write $G \iso H$.
\end{defn}

\begin{leanbox}
\begin{minted}{lean}
structure GroupHom (G H : Type) [Group G] [Group H] where
  toFun : G → H
  map_mul : ∀ a b : G, toFun (a * b) = toFun a * toFun b

-- A homomorphism preserves the identity
theorem GroupHom.map_one [Group G] [Group H] (φ : GroupHom G H) :
    φ.toFun 1 = 1 := by
  have h := φ.map_mul 1 1
  -- h : φ.toFun (1 * 1) = φ.toFun 1 * φ.toFun 1.
  -- Since 1 * 1 = 1 (Group.mul_one), this says
  --   φ.toFun 1 = φ.toFun 1 * φ.toFun 1;
  -- multiplying both sides on the right by (φ.toFun 1)⁻¹ finishes.
  sorry -- Exercise for the reader

-- A homomorphism preserves inverses
theorem GroupHom.map_inv [Group G] [Group H] (φ : GroupHom G H) (a : G) :
    φ.toFun a⁻¹ = (φ.toFun a)⁻¹ := by
  sorry -- Exercise: use map_mul and map_one
\end{minted}
\end{leanbox}

\begin{exercisebox}[title={Exercise 1.4}]
\stepcounter{exercise}
Fill in both \lc{sorry}s above: prove that a group homomorphism preserves the
identity and preserves inverses. \emph{Hint:} For \lc{map\_one}, start with
the equation $\varphi(1) = \varphi(1 \cdot 1) = \varphi(1) \cdot \varphi(1)$
and multiply both sides on the right by $\varphi(1)^{-1}$.
\end{exercisebox}

\begin{exercisebox}[title={Exercise 1.5}]
\stepcounter{exercise}
Define a group homomorphism $\varphi : \Z \to \Z/n\Z$ by
$\varphi(k) = k \bmod n$, and prove \lc{map\_mul} in Lean.
\end{exercisebox}

% ══════════════════════════════════════════════════════════════════════════════
\chapter{Subgroups, Cosets, and Lagrange's Theorem}
% ══════════════════════════════════════════════════════════════════════════════

\epigraph{``It is not knowledge, but the act of learning; not possession, but
the act of getting there, which grants the greatest enjoyment.''}%
{\textit{Carl Friedrich Gauss}}

\section{Subgroups}

A \textbf{subgroup} is a subset of a group that is itself a group under the
same operation. This is the group-theoretic analogue of a ``sub-interface''
or ``inherited class''---it has all the structure of the parent, just on
a smaller set.

\begin{defn}{Subgroup}{subgroup}
Let $G$ be a group. A subset $H \subseteq G$ is a \textbf{subgroup} (written
$H \leq G$) if:
\begin{enumerate}[nosep, label=(\roman*)]
\item \textbf{Identity:} $e \in H$.
\item \textbf{Closure:} If $a, b \in H$, then $a \cdot b \in H$.
\item \textbf{Inverses:} If $a \in H$, then $a^{-1} \in H$.
\end{enumerate}
\end{defn}

\begin{leanbox}
\begin{minted}{lean}
-- Core Lean has no `Set` type (that lives in Mathlib), so we define
-- our own: a set is just a membership predicate.
abbrev Set (α : Type) := α → Prop
instance : Membership α (Set α) where
  mem s a := s a          -- (a ∈ s) unfolds to (s a)

structure Subgroup (G : Type) [Group G] where
  carrier : Set G
  one_mem : (1 : G) ∈ carrier
  mul_mem : ∀ {a b : G}, a ∈ carrier → b ∈ carrier → a * b ∈ carrier
  inv_mem : ∀ {a : G}, a ∈ carrier → a⁻¹ ∈ carrier

-- Membership notation (collection first, element second, matching
-- core Lean's `Membership.mem`)
instance [Group G] : Membership G (Subgroup G) where
  mem H a := a ∈ H.carrier
\end{minted}
\end{leanbox}

\paragraph{Examples.}
\begin{itemize}[nosep]
\item The \textbf{trivial subgroup} $\{e\} \leq G$ and $G \leq G$ itself.
\item The even integers $2\Z \leq \Z$ (under addition).
\item The rotations $\{e, r, r^2\} \leq S_3$ (a subgroup of order 3).
\item For any element $g \in G$, the \textbf{cyclic subgroup}
      $\langle g \rangle = \{g^n : n \in \Z\}$ is a subgroup.
\end{itemize}

\section{Cosets and Lagrange's Theorem}

The key idea: given a subgroup $H \leq G$, we can ``translate'' $H$ by
multiplying all its elements by some fixed group element. The resulting set
is called a \emph{coset}.

\begin{defn}{Coset}{coset}
Let $H \leq G$ and $a \in G$. The \textbf{left coset} of $H$ by $a$ is
\[
  aH \;=\; \{a \cdot h : h \in H\}.
\]
\end{defn}

The crucial property of cosets is that they \emph{partition} $G$: every element
of $G$ belongs to exactly one left coset of $H$. Moreover, all cosets have the
same size as $H$ (because the map $h \mapsto ah$ is a bijection $H \to aH$).

\begin{thm}{Lagrange's Theorem}{lagrange}
Let $G$ be a finite group and $H \leq G$. Then $|H|$ divides $|G|$.
Moreover, the number of distinct left cosets of $H$ in $G$ is
$[G : H] = |G| / |H|$.
\end{thm}

\begin{proof}
The left cosets of $H$ partition $G$ (this requires showing that the coset
relation $a \sim b \iff a^{-1}b \in H$ is an equivalence relation). Each coset
has $|H|$ elements (the map $h \mapsto ah$ is injective). If there are $k$
cosets, then $|G| = k \cdot |H|$, so $|H|$ divides $|G|$.
\end{proof}

\begin{historical}[title={Lagrange and the ``Theorem''}]
Lagrange never stated this theorem in the language of groups---the concept
didn't exist yet. What he showed, in his 1770--1771 \textit{Réflexions sur
la résolution algébrique des équations}, was that certain expressions in the
roots of a polynomial take on a number of distinct values that divides $n!$
(where $n$ is the degree). This is Lagrange's theorem applied to the symmetric
group $S_n$. The abstract formulation had to wait for Galois, Cayley, and
their successors.
\end{historical}

\subsection{Corollaries}

Lagrange's theorem has immediate and powerful consequences:

\begin{cor}{Order of an Element Divides Group Order}{order-divides}
If $G$ is a finite group and $g \in G$, then the order of $g$ (the smallest
positive $n$ such that $g^n = e$) divides $|G|$.
\end{cor}

\begin{proof}
The cyclic subgroup $\langle g \rangle$ has order equal to $\ord(g)$, and
$\langle g \rangle \leq G$, so $\ord(g) \mid |G|$ by Lagrange.
\end{proof}

\begin{cor}{Euler's Theorem}{euler}
If $\gcd(a, n) = 1$, then $a^{\varphi(n)} \equiv 1 \pmod{n}$.
\end{cor}

\begin{proof}
The element $\bar{a} \in (\Z/n\Z)^\times$ has some order $d$ dividing
$|(\Z/n\Z)^\times| = \varphi(n)$ by Corollary~\ref{cor:order-divides}.
So $\varphi(n) = dk$ for some $k$, and
$a^{\varphi(n)} = a^{dk} = (a^d)^k \equiv 1^k = 1$.
\end{proof}

\begin{cor}{Fermat's Little Theorem}{fermat-little}
If $p$ is prime and $p \nmid a$, then $a^{p-1} \equiv 1 \pmod{p}$.
\end{cor}

\begin{proof}
When $p$ is prime, $\varphi(p) = p - 1$, so this is Euler's theorem
specialized to $n = p$.
\end{proof}

\begin{historical}[title={Euler Proves Fermat}]
Fermat stated his ``little theorem'' in 1640 (in a letter to Frénicle de
Bessy) without proof. Euler gave the first published proof in 1736, and
later generalized it to the theorem bearing his own name. The fact that
Fermat's and Euler's theorems are both one-line corollaries of Lagrange's
theorem illustrates the power of the group-theoretic framework: results that
required individual ingenuity become routine consequences of general structure.
\end{historical}

\section{Normal Subgroups and Quotient Groups}

Not every subgroup allows us to form a quotient group. The subgroups that do
are called \emph{normal}.

\begin{defn}{Normal Subgroup}{normal}
A subgroup $N \leq G$ is \textbf{normal} (written $N \normal G$) if
$gNg^{-1} = N$ for all $g \in G$, i.e., conjugation by any element of $G$
maps $N$ to itself.

Equivalently: $aN = Na$ for all $a \in G$ (left cosets equal right cosets).
\end{defn}

\begin{leanbox}
\begin{minted}{lean}
structure NormalSubgroup (G : Type) [Group G] extends Subgroup G where
  conj_mem : ∀ {n : G}, n ∈ carrier → ∀ g : G,
    g * n * g⁻¹ ∈ carrier
\end{minted}
\end{leanbox}

Why normality matters: when $N \normal G$, the set of cosets $G/N = \{aN : a \in G\}$
forms a group under the operation $(aN)(bN) = (ab)N$. This is well-defined
precisely because $N$ is normal. We call $G/N$ the \textbf{quotient group}.

\begin{defn}{Quotient Group}{quotient-group}
For $N \normal G$, the \textbf{quotient group} $G/N$ is the set of left
cosets of $N$, with multiplication $(aN)(bN) = (ab)N$.
\end{defn}

The quotient group $G/N$ has $|G/N| = |G|/|N| = [G:N]$ elements.

\begin{thm}{First Isomorphism Theorem}{first-iso}
If $\varphi : G \to H$ is a group homomorphism, then:
\begin{enumerate}[nosep, label=(\roman*)]
\item $\ker(\varphi) = \{g \in G : \varphi(g) = e_H\}$ is a normal subgroup of $G$.
\item $\im(\varphi) = \{\varphi(g) : g \in G\}$ is a subgroup of $H$.
\item $G / \ker(\varphi) \iso \im(\varphi)$.
\end{enumerate}
\end{thm}

This theorem is the workhorse of group theory---it says that every
homomorphism factors as a projection (onto a quotient) followed by an
injection (into the target).

\begin{exercisebox}[title={Exercise 2.1}]
\stepcounter{exercise}
Compute all subgroups of $S_3$ and determine which are normal.

\emph{Answer (verify in Lean):} $S_3$ has 6 subgroups: $\{e\}$,
$\langle s \rangle$, $\langle sr \rangle$, $\langle sr^2 \rangle$,
$\langle r \rangle = \{e, r, r^2\}$, and $S_3$ itself.
Of these, only $\{e\}$, $\langle r \rangle$, and $S_3$ are normal.
\end{exercisebox}

\begin{exercisebox}[title={Exercise 2.2}]
\stepcounter{exercise}
Prove in Lean that $\ker(\varphi)$ is a normal subgroup for any group
homomorphism $\varphi$.
\end{exercisebox}

% ══════════════════════════════════════════════════════════════════════════════
\chapter{The Subgroup Lattice --- Order Theory Interlude}
% ══════════════════════════════════════════════════════════════════════════════

\epigraph{``In science there is only physics; all the rest is stamp
collecting.''}{\textit{Ernest Rutherford}\\[2pt]
\small (Mathematicians disagree.)}

\section{Why Order Theory?}

In Chapter~2 we computed the subgroups of $S_3$. Now step back and look at the
\emph{structure} of how those subgroups relate to each other. We can order them
by inclusion: $\{e\} \subseteq \langle s \rangle \subseteq S_3$, and so on.
The resulting structure turns out to be extraordinarily important---it is a
\emph{lattice}, and the correspondence at the heart of Galois theory
(Chapter~8) is a bijection between two such lattices.

\begin{historical}[title={Dedekind and the Birth of Lattice Theory}]
Richard Dedekind (1831--1916) was one of the most important algebraists of the
19th century. His contributions include the modern definition of ideals in
rings, the Dedekind cut construction of the real numbers, and---most relevant
for us---the recognition (around 1897) that the subgroups of a group, ordered
by inclusion, form a structure with well-defined ``meet'' and ``join''
operations. He called these structures \textit{Dualgruppen}; the term
``lattice'' was introduced later by Garrett Birkhoff in the 1930s.

Dedekind was also a close collaborator of Cantor and one of the first to
recognize the importance of set-theoretic thinking in algebra. His influence
on the development of abstract algebra is hard to overstate---Emmy Noether,
who revolutionized algebra in the 1920s, explicitly credited Dedekind as her
primary inspiration.
\end{historical}

\section{Partial Orders}

\begin{defn}{Partial Order}{partial-order}
A \textbf{partial order} on a set $P$ is a binary relation $\leq$ that is:
\begin{enumerate}[nosep, label=(\roman*)]
\item \textbf{Reflexive:} $a \leq a$ for all $a \in P$.
\item \textbf{Antisymmetric:} If $a \leq b$ and $b \leq a$, then $a = b$.
\item \textbf{Transitive:} If $a \leq b$ and $b \leq c$, then $a \leq c$.
\end{enumerate}
The pair $(P, \leq)$ is called a \textbf{partially ordered set} (\textbf{poset}).
\end{defn}

\begin{leanbox}
\begin{minted}{lean}
class PartialOrder (P : Type) where
  le : P → P → Prop
  le_refl : ∀ a : P, le a a
  le_antisymm : ∀ a b : P, le a b → le b a → a = b
  le_trans : ∀ a b c : P, le a b → le b c → le a c

instance [PartialOrder P] : LE P where le := PartialOrder.le
\end{minted}
\end{leanbox}

\paragraph{Examples.}
\begin{itemize}[nosep]
\item $(\N, \leq)$ --- the usual ordering on natural numbers.
\item $(\N, \mid)$ --- divisibility. $a \leq b$ means $a \mid b$.
      This is a partial order but not a total order: $2 \nmid 3$ and
      $3 \nmid 2$.
\item The powerset $(\mathcal{P}(S), \subseteq)$ --- subsets ordered by inclusion.
\item The subgroups of $G$, ordered by inclusion.
\end{itemize}

\subsection{Hasse Diagrams}

A \textbf{Hasse diagram} is a visual representation of a poset: elements are
nodes, and we draw an edge from $a$ up to $b$ whenever $a < b$ and there is
no $c$ with $a < c < b$ (a ``covering relation'').

\begin{center}
\begin{tikzpicture}[
    node distance=1.2cm,
    every node/.style={draw, rounded corners, fill=lightblue, minimum width=1.8cm,
    font=\small}
  ]
  \node (S3) {$S_3$};
  \node[below left=1.2cm and 0.2cm of S3] (r) {$\langle r \rangle$};
  \node[below left=1.2cm and -1.5cm of S3] (s) {$\langle s \rangle$};
  \node[below right=1.2cm and -1.5cm of S3] (sr) {$\langle sr \rangle$};
  \node[below right=1.2cm and 0.2cm of S3] (sr2) {$\langle sr^2 \rangle$};
  \node[below=3.5cm of S3] (e) {$\{e\}$};
  \draw[thick] (S3) -- (r);
  \draw[thick] (S3) -- (s);
  \draw[thick] (S3) -- (sr);
  \draw[thick] (S3) -- (sr2);
  \draw[thick] (r) -- (e);
  \draw[thick] (s) -- (e);
  \draw[thick] (sr) -- (e);
  \draw[thick] (sr2) -- (e);
  \node[draw=none, fill=none, below=0.3cm of e] {Subgroup lattice of $S_3$};
\end{tikzpicture}
\end{center}

\section{Lattices}

\begin{defn}{Lattice}{lattice}
A \textbf{lattice} is a poset $(L, \leq)$ in which every pair of elements
$a, b \in L$ has:
\begin{itemize}[nosep]
\item A \textbf{greatest lower bound} (meet, infimum): $a \sqcap b$, satisfying
  $a \sqcap b \leq a$, $a \sqcap b \leq b$, and if $c \leq a$ and $c \leq b$
  then $c \leq a \sqcap b$.
\item A \textbf{least upper bound} (join, supremum): $a \sqcup b$, satisfying
  $a \leq a \sqcup b$, $b \leq a \sqcup b$, and if $a \leq c$ and $b \leq c$
  then $a \sqcup b \leq c$.
\end{itemize}
\end{defn}

\begin{leanbox}
\begin{minted}{lean}
class Lattice (L : Type) extends PartialOrder L where
  meet : L → L → L
  join : L → L → L
  meet_le_left  : ∀ a b : L, le (meet a b) a
  meet_le_right : ∀ a b : L, le (meet a b) b
  le_meet : ∀ a b c : L, le c a → le c b → le c (meet a b)
  le_join_left  : ∀ a b : L, le a (join a b)
  le_join_right : ∀ a b : L, le b (join a b)
  join_le : ∀ a b c : L, le a c → le b c → le (join a b) c

-- (Core Lean has no `Inf`/`Sup` typeclass -- those are Mathlib. We use
-- `Lattice.meet` and `Lattice.join` directly, or add ⊓/⊔ notation.)
\end{minted}
\end{leanbox}

\subsection{The Subgroup Lattice}

\begin{thm}{Subgroups Form a Complete Lattice}{subgroup-lattice}
For any group $G$, the set of subgroups of $G$, ordered by inclusion, forms
a complete lattice:
\begin{itemize}[nosep]
\item $H_1 \sqcap H_2 = H_1 \cap H_2$ (intersection of subgroups is a subgroup).
\item $H_1 \sqcup H_2 = \langle H_1 \cup H_2 \rangle$ (subgroup generated by the union).
\end{itemize}
\end{thm}

\begin{proof}
The intersection of any family of subgroups is a subgroup (check: the identity
is in all of them; if $a, b$ are in all of them, so is $ab$; if $a$ is in all
of them, so is $a^{-1}$). This gives arbitrary meets. The join is defined as
the intersection of all subgroups containing $H_1 \cup H_2$---this exists
because $G$ itself is such a subgroup. Arbitrary meets and a top element give
a complete lattice.
\end{proof}

\begin{warning}
The join of two subgroups is \emph{not} their union. The union of two subgroups
is almost never a subgroup (try $\{0, 2, 4\}$ and $\{0, 3\}$ in $\Z/6\Z$:
their union contains $2$ and $3$ but not $2 + 3 = 5$). The join is the
\emph{smallest subgroup containing} the union.
\end{warning}

\section{Galois Connections}

This is the key concept that will unlock the fundamental theorem of Galois
theory. The definition is abstract, but the payoff is enormous.

\begin{defn}{Galois Connection (antitone)}{galois-connection}
Let $(P, \leq_P)$ and $(Q, \leq_Q)$ be posets. An \textbf{(antitone) Galois
connection} between $P$ and $Q$ is a pair of maps
$\ell : P \to Q$ and $u : Q \to P$ such that:
\[
  q \leq_Q \ell(p) \quad \iff \quad p \leq_P u(q)
  \qquad \text{for all } p \in P, \; q \in Q.
\]
Both $\ell$ and $u$ turn out to be \emph{order-reversing}
(Proposition~\ref{prop:gc-properties}). This is the variant Galois theory
needs: bigger fields $\leftrightarrow$ smaller groups.
\end{defn}

\begin{leanbox}
\begin{minted}{lean}
structure GaloisConnection
    (P Q : Type) [PartialOrder P] [PartialOrder Q] where
  l : P → Q
  u : Q → P
  adjoint : ∀ (p : P) (q : Q), q ≤ l p ↔ p ≤ u q
\end{minted}
\end{leanbox}

\begin{historical}[title={Ore's Abstraction}]
The concept of a Galois connection was abstracted by Oystein Ore in 1944,
who noticed that the specific correspondence Galois discovered between
fields and groups (Chapter~8) is an instance of a general pattern appearing
throughout mathematics. Ore's paper ``Galois Connexions'' extracted the
essential order-theoretic structure and showed it applies far beyond algebra.
The name is fitting: the general concept is named after the specific example
that inspired it.
\end{historical}

\subsection{Properties of Galois Connections}

Every Galois connection automatically satisfies:

\begin{prop}{}{gc-properties}
Let $(\ell, u)$ be an antitone Galois connection between $P$ and $Q$. Then:
\begin{enumerate}[nosep, label=(\roman*)]
\item Both $\ell$ and $u$ are \textbf{order-reversing}:
  if $p_1 \leq_P p_2$ then $\ell(p_2) \leq_Q \ell(p_1)$, and if
  $q_1 \leq_Q q_2$ then $u(q_2) \leq_P u(q_1)$.
\item \textbf{Extensive on $P$:} $p \leq_P u(\ell(p))$ for all $p \in P$.
\item \textbf{Extensive on $Q$:} $q \leq_Q \ell(u(q))$ for all $q \in Q$.
\item \textbf{Idempotence:} $\ell \circ u \circ \ell = \ell$ and
  $u \circ \ell \circ u = u$.
\end{enumerate}
\end{prop}

\begin{proof}
All follow from the adjunction $q \leq_Q \ell(p) \iff p \leq_P u(q)$.
(ii)~Take $q := \ell(p)$: the adjunction reads
$\ell(p) \leq_Q \ell(p) \iff p \leq_P u(\ell(p))$, and the left side holds by
reflexivity. (iii)~Symmetrically, take $p := u(q)$: the adjunction reads
$q \leq_Q \ell(u(q)) \iff u(q) \leq_P u(q)$, and the right side holds by
reflexivity. (i)~If $p_1 \leq_P p_2$, then by (ii) and transitivity
$p_1 \leq_P p_2 \leq_P u(\ell(p_2))$, so $p_1 \leq_P u(\ell(p_2))$; applying the
adjunction with $p := p_1,\ q := \ell(p_2)$ gives $\ell(p_2) \leq_Q \ell(p_1)$.
The argument for $u$ is the mirror image, using (iii). (iv)~follows by combining
(i)--(iii): e.g.\ $\ell$ applied to the unit $p \leq u(\ell(p))$ reverses it to
$\ell(u(\ell(p))) \leq \ell(p)$, while (iii) at $q := \ell(p)$ gives the reverse
inequality; antisymmetry yields $\ell \circ u \circ \ell = \ell$.
\end{proof}

\begin{warning}
There are two conventions in the literature: some authors define Galois
connections with \emph{both maps order-preserving} (the ``monotone'' variant),
and others with \emph{both maps order-reversing} (the ``antitone'' variant, the
one defined above). The Galois correspondence in algebra uses the antitone
variant: bigger fields $\leftrightarrow$ smaller groups. We follow the antitone
convention throughout this book.
\end{warning}

\subsection{Closure Operators}

Every Galois connection gives rise to a \emph{closure operator}: the
composition $u \circ \ell : P \to P$ is extensive ($p \leq u(\ell(p))$),
monotone, and idempotent. The ``closed'' elements---those satisfying
$p = u(\ell(p))$---are exactly the elements in the image of $u$, and
they form a complete lattice.

\begin{defn}{Closure Operator}{closure-op}
A \textbf{closure operator} on a poset $(P, \leq)$ is a map
$\mathrm{cl} : P \to P$ that is:
\begin{enumerate}[nosep, label=(\roman*)]
\item \textbf{Extensive:} $p \leq \mathrm{cl}(p)$.
\item \textbf{Monotone:} $p \leq q \implies \mathrm{cl}(p) \leq \mathrm{cl}(q)$.
\item \textbf{Idempotent:} $\mathrm{cl}(\mathrm{cl}(p)) = \mathrm{cl}(p)$.
\end{enumerate}
\end{defn}

\begin{exercisebox}[title={Exercise 3.1}]
\stepcounter{exercise}
Prove in Lean that $u \circ \ell$ is a closure operator for any (antitone)
Galois connection $(\ell, u)$.
\end{exercisebox}

\begin{exercisebox}[title={Exercise 3.2}]
\stepcounter{exercise}
Define the lattice of divisors of 12, ordered by divisibility. Draw its Hasse
diagram. Verify that it is a lattice by computing $\gcd$ (meet) and $\lcm$
(join) for all pairs.
\end{exercisebox}

\begin{exercisebox}[title={Exercise 3.3}]
\stepcounter{exercise}
Implement \lc{Lattice (Subgroup G)} in Lean for a finite group \lc{G}.
The meet is intersection; the join is the subgroup generated by the union.
Prove the lattice axioms.
\end{exercisebox}

% ══════════════════════════════════════════════════════════════════════════════
\chapter{Rings and Fields --- Adding Multiplication}
% ══════════════════════════════════════════════════════════════════════════════

\epigraph{``In mathematics the art of proposing a question must be held of
higher value than solving it.''}{\textit{Georg Cantor}}

\section{The Rise of Abstract Algebra}

\begin{historical}[title={Hilbert, Noether, and the Abstraction Revolution}]
David Hilbert (1862--1943) and Emmy Noether (1882--1935) transformed algebra
from a collection of computational techniques into a structural science.

Hilbert's 1888 proof of the basis theorem---that every ideal in a polynomial
ring is finitely generated---was purely existential, providing no algorithm for
finding the generators. Paul Gordan reportedly declared: ``This is not
mathematics; this is theology!'' The controversy between Hilbert's abstract
style and Gordan's computational approach foreshadowed the constructivism
debate that continues today.

Emmy Noether took the abstraction further. Her 1921 paper \textit{Idealtheorie
in Ringbereichen} created modern commutative algebra, showing that the right
way to understand a ring is through its \emph{ideals}. As a woman in early
20th-century Germany, she was initially denied a paid position at Göttingen;
Hilbert championed her, declaring to the faculty senate: ``I do not see that
the sex of the candidate is an argument against her admission. After all, we
are a university, not a bathing establishment.''

She was dismissed in 1933 when the Nazis expelled all Jewish faculty, and
emigrated to the United States, where she died at 53. Einstein wrote to the
New York Times that she was ``the most significant creative mathematical
genius thus far produced since the higher education of women began.''
\end{historical}

\section{Ring Axioms}

A ring has two operations: addition (forming an abelian group) and
multiplication (associative, with identity, distributing over addition).

\begin{defn}{Ring}{ring}
A \textbf{ring} $(R, +, \cdot)$ satisfies:
\begin{enumerate}[nosep, label=(\roman*)]
\item $(R, +)$ is an abelian group with identity $0$.
\item Multiplication is associative with identity $1$.
\item \textbf{Distributivity:} $a(b + c) = ab + ac$ and $(a + b)c = ac + bc$.
\end{enumerate}
If $ab = ba$ for all $a, b$, the ring is \textbf{commutative}.
\end{defn}

\begin{leanbox}
\begin{minted}{lean}
class Ring (R : Type) where
  add : R → R → R
  zero : R
  neg : R → R
  mul : R → R → R
  one : R
  add_assoc : ∀ a b c : R, add (add a b) c = add a (add b c)
  add_comm  : ∀ a b : R, add a b = add b a
  zero_add  : ∀ a : R, add zero a = a
  neg_add   : ∀ a : R, add (neg a) a = zero
  mul_assoc : ∀ a b c : R, mul (mul a b) c = mul a (mul b c)
  one_mul   : ∀ a : R, mul one a = a
  mul_one   : ∀ a : R, mul a one = a
  left_distrib  : ∀ a b c : R, mul a (add b c) = add (mul a b) (mul a c)
  right_distrib : ∀ a b c : R, mul (add a b) c = add (mul a c) (mul b c)
\end{minted}
\end{leanbox}

\section{Fields --- Where Division Works}

\begin{defn}{Field}{field}
A \textbf{field} is a commutative ring in which every nonzero element has a
multiplicative inverse: for all $a \neq 0$, there exists $a^{-1}$ with
$a \cdot a^{-1} = 1$.
\end{defn}

\begin{leanbox}
\begin{minted}{lean}
-- Notation instances (just like the Group notation of Chapter 1),
-- so that `a + b`, `a * b`, `0`, `1`, `-a` work on any Ring or Field.
instance [Ring R] : Add R  where add  := Ring.add
instance [Ring R] : Mul R  where mul  := Ring.mul
instance [Ring R] : Zero R where zero := Ring.zero
instance [Ring R] : One R  where one  := Ring.one
instance [Ring R] : Neg R  where neg  := Ring.neg

class Field (F : Type) extends Ring F where
  mul_comm : ∀ a b : F, mul a b = mul b a
  inv : F → F
  mul_inv : ∀ a : F, a ≠ zero → mul a (inv a) = one
  zero_ne_one : zero ≠ one
\end{minted}
\end{leanbox}

\paragraph{Key examples.} $\Q$, $\R$, $\C$ are fields. $\Z$ is a ring but not
a field. $\Z/p\Z$ is a field for prime $p$; $\Z/n\Z$ for composite $n$ has
zero divisors (e.g.\ $2 \cdot 3 = 0$ in $\Z/6\Z$).

\begin{thm}{$\Z/p\Z$ is a Field}{zp-field}
For any prime $p$, $\Z/p\Z$ is a field.
\end{thm}

\begin{proof}
For nonzero $\bar{a}$, since $p$ is prime, $\gcd(a,p) = 1$. Bézout gives
$ax + py = 1$, so $\bar{a}\bar{x} = \bar{1}$ in $\Z/p\Z$.
\end{proof}

\section{Ideals and Quotient Rings}

\begin{defn}{Ideal}{ideal}
An \textbf{ideal} $I \subseteq R$ is an additive subgroup with the absorption
property: $r \cdot a \in I$ and $a \cdot r \in I$ for all $r \in R$, $a \in I$.
\end{defn}

Ideals are to rings what normal subgroups are to groups: the quotient
$R/I$ inherits a ring structure precisely when $I$ is an ideal. The connection
to fields:

\begin{thm}{}{quotient-field}
$R/I$ is a field if and only if $I$ is a \textbf{maximal ideal}.
\end{thm}

This will be the key to constructing field extensions in Chapter~6.

\begin{exercisebox}[title={Exercise 4.1}]
\stepcounter{exercise}
Implement \lc{Field (Fin 5)} in Lean. Compute all multiplicative inverses
in $\Z/5\Z$.
\end{exercisebox}

\begin{exercisebox}[title={Exercise 4.2}]
\stepcounter{exercise}
Show $\Z/6\Z$ is not a field: find $a, b$ with $a \neq 0$, $b \neq 0$,
$ab = 0$.
\end{exercisebox}

% ══════════════════════════════════════════════════════════════════════════════
\chapter{Polynomials Over a Field}
% ══════════════════════════════════════════════════════════════════════════════

\epigraph{``I have had my results for a long time, but I do not yet know how I
am to arrive at them.''}{\textit{Carl Friedrich Gauss}}

\section{The 250-Year Quest}

\begin{historical}[title={The Hunt for the Quintic Formula}]
The Babylonians solved quadratics around 2000~BCE. The cubic formula was found
around 1515--1535 and published (amid accusations of oath-breaking) by Cardano
in 1545. Ferrari found the quartic formula shortly after. Then: 250 years of
silence. Euler, Lagrange, Bézout, and many others attempted the quintic, all
unsuccessfully. Lagrange (1770) analyzed why the lower-degree methods worked
and concluded, presciently, that the same approach couldn't work for degree~5%
---but he couldn't prove impossibility. The resolution came from Abel (1824)
and Galois (1832), whose story we tell in Chapter~10.
\end{historical}

\section{The Polynomial Ring $F[x]$}

\begin{defn}{Polynomial Ring}{poly-ring}
$F[x]$ is the set of formal expressions
$f(x) = a_n x^n + \cdots + a_1 x + a_0$ with $a_i \in F$. The
\textbf{degree} is $\deg(f) = n$ (where $a_n \neq 0$).
\end{defn}

\begin{leanbox}
\begin{minted}{lean}
-- Polynomials as coefficient lists (least-significant first)
structure Polynomial (F : Type) [Field F] where
  coeffs : List F
  deriving Repr

namespace Polynomial

def eval [Field F] (p : Polynomial F) (x : F) : F :=
  p.coeffs.foldr (fun c acc => Ring.add (Ring.mul acc x) c) Ring.zero

def degree [Field F] (p : Polynomial F) : Nat :=
  if p.coeffs.length = 0 then 0 else p.coeffs.length - 1

end Polynomial
\end{minted}
\end{leanbox}

\section{Division and Irreducibility}

\begin{thm}{Division Algorithm}{poly-division}
For $f, g \in F[x]$ with $g \neq 0$, there exist unique $q, r \in F[x]$
such that $f = gq + r$ and $\deg(r) < \deg(g)$.
\end{thm}

\begin{defn}{Irreducible Polynomial}{irreducible}
$f \in F[x]$ of degree $\geq 1$ is \textbf{irreducible} over $F$ if it
cannot be factored as $f = gh$ with $\deg(g), \deg(h) \geq 1$.
\end{defn}

\begin{warning}
Irreducibility depends on the field! $x^2 + 1$ is irreducible over $\R$
but factors as $(x+i)(x-i)$ over $\C$, and as $(x+1)^2$ over $\Z/2\Z$.
\end{warning}

\begin{thm}{Unique Factorization}{poly-uft}
Every nonzero $f \in F[x]$ factors uniquely (up to order and units) into
irreducible polynomials.
\end{thm}

\begin{exercisebox}[title={Exercise 5.1}]
\stepcounter{exercise}
Implement polynomial long division in Lean.
\end{exercisebox}

\begin{exercisebox}[title={Exercise 5.2}]
\stepcounter{exercise}
Prove that $x^3 - 2$ is irreducible over $\Q$ using the rational root theorem.
\end{exercisebox}

% ══════════════════════════════════════════════════════════════════════════════
\chapter{Field Extensions}
% ══════════════════════════════════════════════════════════════════════════════

\epigraph{``God made the integers; all else is the work of man.''}%
{\textit{Leopold Kronecker}}

\section{Extending the Number System}

\begin{historical}[title={Kronecker vs.\ Dedekind}]
Leopold Kronecker (1823--1891) insisted that new objects must be
\emph{constructed} explicitly---a proto-constructivist. His approach to field
extensions was concrete: to adjoin a root of an irreducible polynomial $p(x)$,
work in $F[x]/(p(x))$. Dedekind preferred an axiomatic approach, defining
fields by their properties and proving existence abstractly. Both give the
same answer; in Lean, we naturally follow Kronecker.
\end{historical}

\section{The Basic Construction}

If $p(x) \in F[x]$ is irreducible, we ``add a root'' by quotienting:
$(p(x))$ is maximal, so $F[x]/(p(x))$ is a field.

\begin{defn}{Simple Algebraic Extension}{simple-ext}
For irreducible $p \in F[x]$ of degree $n$:
\[
  F(\alpha) = F[x]/(p(x)) \iso
  \{a_0 + a_1\alpha + \cdots + a_{n-1}\alpha^{n-1} : a_i \in F\}
\]
where $\alpha$ is the image of $x$, satisfying $p(\alpha) = 0$.
\end{defn}

\begin{leanbox}
\begin{minted}{lean}
-- ℚ(√2): elements are a + b√2 with a, b ∈ ℚ
structure QSqrt2 where
  a : Rat
  b : Rat
  deriving Repr, DecidableEq

namespace QSqrt2

def add (x y : QSqrt2) : QSqrt2 := ⟨x.a + y.a, x.b + y.b⟩

def mul (x y : QSqrt2) : QSqrt2 :=
  ⟨x.a * y.a + 2 * x.b * y.b, x.a * y.b + y.a * x.b⟩

-- (a + b√2)⁻¹ = (a - b√2)/(a² - 2b²)
def inv (x : QSqrt2) : QSqrt2 :=
  let d := x.a * x.a - 2 * x.b * x.b
  ⟨x.a / d, -x.b / d⟩

def sqrt2 : QSqrt2 := ⟨0, 1⟩

-- Verify: √2 · √2 = 2
example : mul sqrt2 sqrt2 = ⟨2, 0⟩ := by native_decide

end QSqrt2
\end{minted}
\end{leanbox}

\section{Degree of an Extension}

\begin{defn}{Degree}{degree-ext}
$[E:F]$ is the dimension of $E$ as an $F$-vector space.
\end{defn}

For $F(\alpha)$ with $\alpha$ a root of irreducible $p$ of degree $n$:
$[F(\alpha):F] = n$, with basis $\{1, \alpha, \ldots, \alpha^{n-1}\}$.

\begin{thm}{Tower Law}{tower-law}
For $F \subseteq K \subseteq L$: $[L:F] = [L:K] \cdot [K:F]$.
\end{thm}

\section{Splitting Fields}

\begin{defn}{Splitting Field}{splitting-field}
The \textbf{splitting field} of $f \in F[x]$ is the smallest extension of
$F$ in which $f$ factors into linear factors.
\end{defn}

\begin{thm}{Existence and Uniqueness}{splitting-exists}
Every nonconstant polynomial has a splitting field, unique up to isomorphism.
\end{thm}

\begin{historical}[title={The Fundamental Theorem of Algebra}]
Every polynomial over $\C$ splits completely---$\C$ is algebraically closed.
Gauss gave the first rigorous proof in his 1799 doctoral dissertation and
returned to this theorem throughout his life, giving four different proofs.
The fact that the ``fundamental theorem of algebra'' requires analysis (not
just algebra) for its proof is one of the great ironies of mathematics.
\end{historical}

\begin{exercisebox}[title={Exercise 6.1}]
\stepcounter{exercise}
Construct $\Q(\sqrt[3]{2})$ in Lean: elements $a + b\sqrt[3]{2} + c(\sqrt[3]{2})^2$. Implement arithmetic using $(\sqrt[3]{2})^3 = 2$.
\end{exercisebox}

\begin{exercisebox}[title={Exercise 6.2}]
\stepcounter{exercise}
Show $[\Q(\sqrt{2}, \sqrt{3}):\Q] = 4$ using the tower law.
\end{exercisebox}

% ══════════════════════════════════════════════════════════════════════════════
\part{The Galois Correspondence}
% ══════════════════════════════════════════════════════════════════════════════

% ══════════════════════════════════════════════════════════════════════════════
\chapter{Field Automorphisms and the Galois Group}
% ══════════════════════════════════════════════════════════════════════════════

\epigraph{``Je n'ai pas le temps.''\\ \normalfont(I have no time.)}%
{\textit{Évariste Galois}, writing in the margins\\the night before his duel}

\section{Galois's Story}

\begin{historical}[title={Évariste Galois (1811--1832)}]
The life of Évariste Galois is the most dramatic story in mathematics. Born in
Bourg-la-Reine, near Paris, he was recognized as a mathematical prodigy by his
teachers at the Lycée Louis-le-Grand. At 17, he submitted his first paper to
the French Academy of Sciences---it was lost by Cauchy. He submitted again to
the Academy's prize competition; the paper was sent to Fourier for review, but
Fourier died before reading it, and the manuscript was never found. He was
rejected twice from the École Polytechnique (the premier mathematics school in
France), apparently because his answers were too far ahead of his examiners.

Meanwhile, Paris was in political turmoil. Galois was a passionate republican,
arrested twice for political activities. During his second imprisonment in
1832, he fell in love with Stéphanie-Félicie Poterin du Motel---the details
are murky, but it appears a duel resulted, possibly engineered by political
enemies. The night before the duel, on May 29, 1832, Galois wrote a long
letter to his friend Auguste Chevalier, frantically summarizing his
mathematical discoveries. He scrawled ``I have no time'' in the margins.

The next morning, Galois was shot in the abdomen. He died the following day,
at the age of twenty. His brother Alfred gathered his manuscripts. Liouville
finally published them in 1846, fourteen years later, recognizing their
revolutionary importance.

What Galois created in those few feverish years was nothing less than a new
way of thinking about algebra: instead of asking ``can I solve this equation?''
he asked ``what are the \emph{symmetries} of this equation's solutions?'' The
symmetry group---now called the Galois group---controls everything.
\end{historical}

\section{Field Automorphisms}

\begin{defn}{Field Automorphism}{field-aut}
Let $E/F$ be a field extension. A \textbf{field automorphism} of $E$ fixing
$F$ is a bijection $\sigma : E \to E$ such that:
\begin{enumerate}[nosep, label=(\roman*)]
\item $\sigma(a + b) = \sigma(a) + \sigma(b)$ for all $a, b \in E$.
\item $\sigma(a \cdot b) = \sigma(a) \cdot \sigma(b)$ for all $a, b \in E$.
\item $\sigma(a) = a$ for all $a \in F$.
\end{enumerate}
\end{defn}

\begin{leanbox}
\begin{minted}{lean}
structure FieldAut (E F : Type) [Field E] [Field F] where
  toFun  : E → E
  invFun : E → E
  embed  : F → E            -- the inclusion map F ↪ E
  left_inv  : ∀ x, invFun (toFun x) = x
  right_inv : ∀ x, toFun (invFun x) = x
  map_add : ∀ a b, toFun (a + b) = toFun a + toFun b
  map_mul : ∀ a b, toFun (a * b) = toFun a * toFun b
  fixes_base : ∀ a : F, toFun (embed a) = embed a
\end{minted}
\end{leanbox}

\begin{defn}{Galois Group}{galois-group}
The \textbf{Galois group} of $E/F$, written $\Gal(E/F)$, is the set of all
field automorphisms of $E$ fixing $F$, with composition as the group operation.
\end{defn}

\section{First Example: $\Gal(\Q(\sqrt{2})/\Q)$}

What automorphisms of $\Q(\sqrt{2})$ fix $\Q$? An automorphism $\sigma$ is
determined by $\sigma(\sqrt{2})$, since every element is $a + b\sqrt{2}$ and
$\sigma(a + b\sqrt{2}) = a + b \cdot \sigma(\sqrt{2})$.

Now $\sigma(\sqrt{2})$ must be a root of $x^2 - 2$ (since $\sigma$ preserves
polynomial relations). The roots are $\sqrt{2}$ and $-\sqrt{2}$. So there are
exactly two automorphisms:
\begin{align*}
  \id &: a + b\sqrt{2} \mapsto a + b\sqrt{2} \\
  \sigma &: a + b\sqrt{2} \mapsto a - b\sqrt{2}
\end{align*}

Thus $\Gal(\Q(\sqrt{2})/\Q) = \{\id, \sigma\} \iso \Z/2\Z$.

\begin{leanbox}
\begin{minted}{lean}
-- The nontrivial automorphism of ℚ(√2)
def conjugate : QSqrt2 → QSqrt2
  | ⟨a, b⟩ => ⟨a, -b⟩

-- Verify: conjugate preserves multiplication.
-- (Mathlib's `ring` tactic would close this in one line; in core Lean
--  we discharge the two ℚ-components with the ring lemmas by hand.)
theorem conjugate_mul (x y : QSqrt2) :
    conjugate (QSqrt2.mul x y) = QSqrt2.mul (conjugate x) (conjugate y) := by
  simp only [conjugate, QSqrt2.mul, QSqrt2.mk.injEq]
  refine ⟨?_, ?_⟩
  · show x.a * y.a + 2 * x.b * y.b = x.a * y.a + 2 * -x.b * -y.b
    have h1 : (2 : Rat) * -x.b = -(2 * x.b) := by rw [Rat.mul_neg]
    have h2 : -(2 * x.b) * -y.b = 2 * x.b * y.b := by
      rw [Rat.neg_mul, Rat.mul_neg, Rat.neg_neg]
    rw [h1, h2]
  · show -(x.a * y.b + y.a * x.b) = x.a * -y.b + y.a * -x.b
    rw [Rat.mul_neg, Rat.mul_neg, Rat.neg_add]

-- Verify: conjugate is an involution (applying it twice gives identity)
theorem conjugate_conjugate (x : QSqrt2) :
    conjugate (conjugate x) = x := by
  simp [conjugate]

-- Verify: conjugate fixes ℚ (elements with b = 0)
theorem conjugate_fixes_rat (a : Rat) :
    conjugate ⟨a, 0⟩ = ⟨a, 0⟩ := by
  simp [conjugate]
\end{minted}
\end{leanbox}

\section{Normal and Separable Extensions}

\begin{defn}{Normal Extension}{normal-ext}
$E/F$ is \textbf{normal} if every irreducible polynomial in $F[x]$ that has
at least one root in $E$ splits completely in $E$.
\end{defn}

\begin{defn}{Separable Extension}{separable-ext}
$E/F$ is \textbf{separable} if every element of $E$ is a root of a
\emph{separable} polynomial over $F$ (one with no repeated roots).
\end{defn}

\begin{intuition}
\textbf{Why haven't we worried about separability before?} Because every field
in this book so far is $\Q$ or a subfield of $\C$, and \emph{over any field of
characteristic $0$ separability is automatic}: an irreducible polynomial and its
formal derivative can never share a root there, so repeated roots are
impossible. The hypothesis only starts to bite in characteristic $p$---which is
exactly why it reappears in the finite-field chapter (Chapter~11), where
$\GF(p^n)$ lives in characteristic $p$. For our number-field examples you may
read ``Galois'' as simply ``normal.''
\end{intuition}

\begin{defn}{Galois Extension}{galois-ext}
$E/F$ is a \textbf{Galois extension} if it is both normal and separable.
Equivalently, $|\Gal(E/F)| = [E:F]$.
\end{defn}

\begin{historical}[title={Dedekind and Artin}]
Galois himself didn't separate the concepts of normality and separability---the
distinction was clarified by Dedekind in the 1870s and given its modern form by
Emil Artin in the 1940s. Artin's reformulation of Galois theory, replacing
Galois's original (somewhat obscure) arguments with clean linear algebra, is
the version we follow in this book.
\end{historical}

\begin{intuition}
Normal ensures there are \emph{enough} automorphisms. Separable ensures
they're all \emph{distinct}. Together: $|\Gal(E/F)| = [E:F]$, so the Galois
group is ``as large as it could possibly be.''
\end{intuition}

\begin{exercisebox}[title={Exercise 7.1}]
\stepcounter{exercise}
Compute $\Gal(\Q(\sqrt{2}, \sqrt{3})/\Q)$. Show it is isomorphic to
$\Z/2\Z \times \Z/2\Z$ (the Klein four-group) by exhibiting all four
automorphisms.
\end{exercisebox}

\begin{exercisebox}[title={Exercise 7.2}]
\stepcounter{exercise}
Show that $\Q(\sqrt[3]{2})/\Q$ is \emph{not} a normal extension.
\emph{Hint:} $x^3 - 2$ has one root in $\Q(\sqrt[3]{2})$ but not all three
(the other two are complex).
\end{exercisebox}

% ══════════════════════════════════════════════════════════════════════════════
\chapter{The Fundamental Theorem of Galois Theory}
% ══════════════════════════════════════════════════════════════════════════════

\epigraph{``Mathematics is the queen of the sciences, and number theory is
the queen of mathematics.''}{\textit{Carl Friedrich Gauss}}

\section{The Main Event}

This is the theorem that Galois glimpsed in 1832, Dedekind clarified in the
1870s, and Artin perfected in the 1940s. It connects two seemingly unrelated
structures---intermediate fields and subgroups---through an order-reversing
bijection.

\begin{thm}{Fundamental Theorem of Galois Theory}{ftgt}
Let $E/F$ be a finite Galois extension with Galois group
$G = \Gal(E/F)$. There is an order-reversing bijection
\[
\begin{tikzcd}[row sep=0.5cm]
  \{\text{intermediate fields } K : F \subseteq K \subseteq E\}
  \arrow[r, shift left, "\fixgrp"]
  &
  \{\text{subgroups } H \leq G\}
  \arrow[l, shift left, "\fixfld"]
\end{tikzcd}
\]
given by the \textbf{fixing-group} map (field $\mapsto$ group) and the
\textbf{fixed-field} map (group $\mapsto$ field):
\begin{align*}
  \fixgrp(K) &= \Gal(E/K) = \{\sigma \in G : \sigma(x) = x \text{ for all } x \in K\} \\
  \fixfld(H) &= E^H = \{x \in E : \sigma(x) = x \text{ for all } \sigma \in H\}
\end{align*}
(These are the standard names, and the ones the Lean code below uses:
\lc{fixingGroup} and \lc{fixedField}.) Moreover:
\begin{enumerate}[nosep, label=(\roman*)]
\item $\fixgrp$ and $\fixfld$ are mutual inverses: $\fixfld(\fixgrp(K)) = K$ and
  $\fixgrp(\fixfld(H)) = H$.
\item The correspondence \textbf{reverses order}: $K_1 \subseteq K_2$
  iff $\fixgrp(K_2) \leq \fixgrp(K_1)$.
\item \textbf{Degrees match indices:} $[E:K] = |\fixgrp(K)|$ and
  $[K:F] = [G : \fixgrp(K)]$.
\item $K/F$ is a normal (Galois) extension iff $\fixgrp(K) \normal G$.
  In that case, $\Gal(K/F) \iso G / \fixgrp(K)$.
\end{enumerate}
\end{thm}

\begin{warning}
This is the deepest theorem in the book, and---unlike everything up to
Chapter~7---we \emph{state} it rather than prove it in Lean. A from-scratch,
core-Lean proof (building intermediate fields, the Galois group, and Artin's
linear-algebra argument) is a substantial project well beyond our scope; the
Lean block below is an honest \emph{skeleton} with \lc{sorry} in place of the
hard proofs. Take the theorem on faith; the worked example in Chapter~9 will
make it concrete.
\end{warning}

\begin{intuition}
The fundamental theorem says: \textbf{the structure of a field extension is
entirely determined by the structure of its symmetry group.} Bigger
intermediate fields have \emph{fewer} symmetries (more elements are pinned
down), so the corresponding subgroup is smaller. This is an \emph{antitone
Galois connection}---exactly the concept from Chapter~3!

In fact, this \emph{is} the original Galois connection. The general concept
(Chapter~3) was abstracted from this specific example by Ore in 1944.
\end{intuition}

\begin{leanbox}
\begin{minted}{lean}
-- ILLUSTRATIVE ONLY: this block sketches the shape of the theorem and
-- does NOT type-check as printed. It references names we have not built
-- from scratch (`Gal E F`, `IntermediateField`, set-builder `{ x | … }`),
-- and every definition ends in `sorry`. Read it as pseudocode.

-- The fixed-field map: subgroup H ↦ elements of E fixed by every σ ∈ H
def fixedField (H : Subgroup (Gal E F)) : IntermediateField E F where
  carrier := { x : E | ∀ σ ∈ H.carrier, σ.toFun x = x }
  sorry  -- proof obligation: this carrier is a subfield

-- The fixing-group map: intermediate field K ↦ automorphisms fixing K
def fixingGroup (K : IntermediateField E F) : Subgroup (Gal E F) where
  carrier := { σ : Gal E F | ∀ x ∈ K.carrier, σ.toFun x = x }
  sorry  -- proof obligation: this carrier is a subgroup

-- These maps form an ANTITONE Galois connection between the lattice of
-- intermediate fields and the lattice of subgroups (Definition 3.x):
theorem galois_connection :
    ∀ K H, H ≤ fixingGroup K ↔ K ≤ fixedField H := by
  sorry  -- the core of the fundamental theorem
\end{minted}
\end{leanbox}

\section{The Two Lattices}

The fundamental theorem is a bijection between two lattices:

\begin{center}
\begin{tikzpicture}[
    field/.style={draw, rounded corners, fill=lightblue, minimum width=2cm,
      font=\small, align=center},
    grp/.style={draw, rounded corners, fill=lightgreen, minimum width=2cm,
      font=\small, align=center},
    node distance=1.3cm,
  ]
  % Fields side
  \node[field] (E) {$E$};
  \node[field, below left=1.3cm and 0.3cm of E] (K1) {$K_1$};
  \node[field, below right=1.3cm and 0.3cm of E] (K2) {$K_2$};
  \node[field, below=3cm of E] (F) {$F$};
  \draw[thick] (E) -- (K1);
  \draw[thick] (E) -- (K2);
  \draw[thick] (K1) -- (F);
  \draw[thick] (K2) -- (F);
  \node[above=0.3cm of E, font=\bfseries\color{deepblue}] {Intermediate Fields};

  % Arrow
  \draw[<->, very thick, accentpurple] (3.5, -1.5) -- node[above, font=\scriptsize] {$\fixgrp / \fixfld$} (5.5, -1.5);

  % Groups side (REVERSED order!)
  \node[grp] at (9, 0) (trivial) {$\{e\}$};
  \node[grp, below left=1.3cm and 0.3cm of trivial] (H2) {$H_2$};
  \node[grp, below right=1.3cm and 0.3cm of trivial] (H1) {$H_1$};
  \node[grp, below=3cm of trivial] (G) {$G$};
  \draw[thick] (trivial) -- (H2);
  \draw[thick] (trivial) -- (H1);
  \draw[thick] (H2) -- (G);
  \draw[thick] (H1) -- (G);
  \node[above=0.3cm of trivial, font=\bfseries\color{accentgreen}] {Subgroups of $G$};
\end{tikzpicture}
\end{center}

Notice the \textbf{reversal}: $E$ (the biggest field) corresponds to $\{e\}$
(the smallest subgroup), and $F$ (the smallest field) corresponds to $G$ (the
biggest subgroup).

\begin{exercisebox}[title={Exercise 8.1}]
\stepcounter{exercise}
Verify the fundamental theorem for $\Q(\sqrt{2}, \sqrt{3})/\Q$:
list all intermediate fields, list all subgroups of
$\Gal(\Q(\sqrt{2}, \sqrt{3})/\Q) \iso \Z/2\Z \times \Z/2\Z$,
and check the correspondence.
\end{exercisebox}

% ══════════════════════════════════════════════════════════════════════════════
\chapter{Worked Example --- The Splitting Field of $x^4 - 2$}
% ══════════════════════════════════════════════════════════════════════════════

\section{Building the Splitting Field}

The roots of $x^4 - 2$ are $\sqrt[4]{2}$, $i\sqrt[4]{2}$,
$-\sqrt[4]{2}$, and $-i\sqrt[4]{2}$. We need both $\sqrt[4]{2}$ and $i$
to express all roots, so the splitting field is $E = \Q(\sqrt[4]{2}, i)$.

By the tower law:
\[
  [E:\Q] = [E:\Q(\sqrt[4]{2})] \cdot [\Q(\sqrt[4]{2}):\Q] = 2 \cdot 4 = 8.
\]

So $|\Gal(E/\Q)| = 8$, and the Galois group has 8 elements.

\section{Computing the Galois Group}

An automorphism $\sigma \in \Gal(E/\Q)$ is determined by:
\begin{itemize}[nosep]
\item $\sigma(\sqrt[4]{2})$: must be a root of $x^4 - 2$, so one of
  $\{\sqrt[4]{2},\; i\sqrt[4]{2},\; -\sqrt[4]{2},\; -i\sqrt[4]{2}\}$
  (4 choices).
\item $\sigma(i)$: must be a root of $x^2 + 1$, so either $i$ or $-i$
  (2 choices).
\end{itemize}
Total: $4 \times 2 = 8$ automorphisms, matching $[E:\Q] = 8$. \checkmark

Let $\rho$ be the automorphism with $\rho(\sqrt[4]{2}) = i\sqrt[4]{2}$,
$\rho(i) = i$ (a ``rotation''), and $\tau$ the one with
$\tau(\sqrt[4]{2}) = \sqrt[4]{2}$, $\tau(i) = -i$ (``complex conjugation'').
Then $\rho$ has order 4 and $\tau$ has order 2, and they satisfy
$\tau\rho\tau = \rho^{-1}$. This is exactly the presentation of the
\textbf{dihedral group} $D_4$---the symmetry group of a square!

\[
  \Gal(\Q(\sqrt[4]{2}, i)/\Q) \;\iso\; D_4
  \;=\; \langle \rho, \tau \mid \rho^4 = \tau^2 = e,\;
  \tau\rho\tau = \rho^{-1} \rangle
\]

\section{The Correspondence}

$D_4$ has exactly 10 subgroups ($1$ trivial, $5$ of order~2, $3$ of order~4, and
$D_4$ itself), so by the fundamental theorem there are exactly 10 intermediate
fields. Here are the two lattices, side by side. The subgroup lattice is drawn
\emph{upside down} (with $\{e\}$ at the top) so that corresponding nodes sit at
the same height: the correspondence \emph{reverses} inclusion, so $E$ pairs with
$\{e\}$ and $\Q$ pairs with $D_4$.

\begin{center}
\begin{tikzpicture}[
    fld/.style={draw, rounded corners, fill=lightblue!80, font=\scriptsize,
      inner sep=2pt},
    grp/.style={draw, rounded corners, fill=lightgreen!80, font=\scriptsize,
      inner sep=2pt},
    scale=0.82, transform shape,
  ]
  % ── FIELDS: Q at bottom, E at top ──────────────────────────────
  \node[fld] (QQ)  at (0,    0) {$\Q$};
  % degree 2 over Q
  \node[fld] (Qs2) at (-2,   2) {$\Q(\sqrt{2})$};
  \node[fld] (Qi)  at (0,    2) {$\Q(i)$};
  \node[fld] (Qis) at (2,    2) {$\Q(i\sqrt{2})$};
  % degree 4 over Q
  \node[fld] (Q4)  at (-3.4, 4) {$\Q(\sqrt[4]{2})$};
  \node[fld] (Qi4) at (-1.7, 4) {$\Q(i\sqrt[4]{2})$};
  \node[fld] (Qsi) at (0,    4) {$\Q(\sqrt{2},i)$};
  \node[fld] (Qp)  at (1.7,  4) {$\Q((1{+}i)\sqrt[4]{2})$};
  \node[fld] (Qm)  at (3.4,  4) {$\Q((1{-}i)\sqrt[4]{2})$};
  % top
  \node[fld] (E)   at (0,    6) {$E=\Q(\sqrt[4]{2},i)$};

  \draw (QQ) -- (Qs2); \draw (QQ) -- (Qi); \draw (QQ) -- (Qis);
  \draw (Qs2) -- (Q4); \draw (Qs2) -- (Qi4); \draw (Qs2) -- (Qsi);
  \draw (Qi)  -- (Qsi);
  \draw (Qis) -- (Qsi); \draw (Qis) -- (Qp); \draw (Qis) -- (Qm);
  \draw (Q4) -- (E); \draw (Qi4) -- (E); \draw (Qsi) -- (E);
  \draw (Qp) -- (E); \draw (Qm) -- (E);
  \node[font=\scriptsize\bfseries\color{deepblue}] at (0,6.9) {Intermediate fields};

  % correspondence arrow
  \node at (6.5, 3) {\large $\longleftrightarrow$};

  % ── SUBGROUPS: D4 at bottom, {e} at top (upside down) ──────────
  \begin{scope}[xshift=13cm]
  \node[grp] (D4)  at (0,    0) {$D_4$};
  % order 4  (pair with degree-2 fields)
  \node[grp] (r2t) at (-2,   2) {$\langle\rho^2,\tau\rangle$};
  \node[grp] (r)   at (0,    2) {$\langle\rho\rangle$};
  \node[grp] (r2rt)at (2,    2) {$\langle\rho^2,\rho\tau\rangle$};
  % order 2  (pair with degree-4 fields)
  \node[grp] (t)   at (-3.4, 4) {$\langle\tau\rangle$};
  \node[grp] (r2ta)at (-1.7, 4) {$\langle\rho^2\tau\rangle$};
  \node[grp] (r2)  at (0,    4) {$\langle\rho^2\rangle$};
  \node[grp] (rt)  at (1.7,  4) {$\langle\rho\tau\rangle$};
  \node[grp] (r3t) at (3.4,  4) {$\langle\rho^3\tau\rangle$};
  % top
  \node[grp] (e1)  at (0,    6) {$\{e\}$};

  \draw (D4) -- (r2t); \draw (D4) -- (r); \draw (D4) -- (r2rt);
  \draw (r2t) -- (t); \draw (r2t) -- (r2ta); \draw (r2t) -- (r2);
  \draw (r)   -- (r2);
  \draw (r2rt) -- (r2); \draw (r2rt) -- (rt); \draw (r2rt) -- (r3t);
  \draw (t) -- (e1); \draw (r2ta) -- (e1); \draw (r2) -- (e1);
  \draw (rt) -- (e1); \draw (r3t) -- (e1);
  \node[font=\scriptsize\bfseries\color{accentgreen}] at (0,6.9) {Subgroups of $D_4$};
  \end{scope}
\end{tikzpicture}
\end{center}

\noindent The vertically-aligned pairs are the correspondence:
$E\leftrightarrow\{e\}$;
$\Q(\sqrt[4]{2})\leftrightarrow\langle\tau\rangle$,
$\Q(i\sqrt[4]{2})\leftrightarrow\langle\rho^2\tau\rangle$,
$\Q(\sqrt{2},i)\leftrightarrow\langle\rho^2\rangle$,
$\Q((1{+}i)\sqrt[4]{2})\leftrightarrow\langle\rho\tau\rangle$,
$\Q((1{-}i)\sqrt[4]{2})\leftrightarrow\langle\rho^3\tau\rangle$;
$\Q(\sqrt{2})\leftrightarrow\langle\rho^2,\tau\rangle$,
$\Q(i)\leftrightarrow\langle\rho\rangle$,
$\Q(i\sqrt{2})\leftrightarrow\langle\rho^2,\rho\tau\rangle$;
$\Q\leftrightarrow D_4$.

\begin{exercisebox}[title={Exercise 9.1}]
\stepcounter{exercise}
Verify the correspondence for at least three intermediate field--subgroup
pairs: show that the automorphisms fixing the field form the claimed subgroup.
\end{exercisebox}

% ══════════════════════════════════════════════════════════════════════════════
\chapter{Solvability --- Why There Is No Quintic Formula}
% ══════════════════════════════════════════════════════════════════════════════

\epigraph{``Abel has left mathematicians enough to keep them busy for five
hundred years.''}{\textit{Charles Hermite}}

\section{Two Tragic Heroes}

\begin{historical}[title={Abel and Galois}]
Niels Henrik Abel (1802--1829) grew up in poverty in Norway. At 22, he proved
that the general quintic equation cannot be solved by radicals---the first
impossibility result in algebra. He sent his proof to Gauss (who ignored it)
and to Cauchy at the French Academy (who lost the manuscript). Abel spent two
years traveling Europe trying to gain recognition, returned to Norway in debt,
contracted tuberculosis, and died at 26. Two days later, a letter arrived
offering him a prestigious professorship in Berlin.

Galois (Chapter~7) went further: he found the exact criterion for when a
polynomial \emph{is} solvable by radicals. The criterion is stated in terms of
the Galois group.

Both men died tragically young---Abel at 26, Galois at 20. Their work was
recognized only posthumously. Today, abelian groups bear Abel's name, and
Galois theory is the cornerstone of modern algebra.
\end{historical}

\section{Solvable Groups}

\begin{defn}{Solvable Group}{solvable}
A group $G$ is \textbf{solvable} if there exists a chain of subgroups
\[
  \{e\} = G_0 \normal G_1 \normal G_2 \normal \cdots \normal G_n = G
\]
where each quotient $G_{i+1}/G_i$ is abelian.
\end{defn}

\begin{leanbox}
\begin{minted}{lean}
-- ILLUSTRATIVE ONLY (does not type-check): `IsNormal`, `IsAbelian`,
-- and a two-argument `Quotient` are not defined here -- this sketches
-- the definition rather than building it from scratch.
-- A composition series: chain of normal subgroups with abelian quotients
def IsSolvable (G : Type) [Group G] : Prop :=
  ∃ (n : Nat) (chain : Fin (n + 1) → Subgroup G),
    chain ⟨0, by omega⟩ = ⊥ ∧              -- starts at {e}
    chain ⟨n, by omega⟩ = ⊤ ∧              -- ends at G
    ∀ i : Fin n,                             -- each step is normal with
      IsNormal (chain i) (chain i.succ) ∧    -- abelian quotient
      IsAbelian (Quotient (chain i.succ) (chain i))
\end{minted}
\end{leanbox}

\section{The Solvability Criterion}

\begin{thm}{Galois's Criterion}{galois-criterion}
A polynomial $f \in F[x]$ is solvable by radicals if and only if its
Galois group $\Gal(E/F)$ (where $E$ is the splitting field of $f$) is a
solvable group.
\end{thm}

\begin{warning}
Like the fundamental theorem (Chapter~8), Galois's criterion is a deep result
we \emph{state and use} but do not formalize from scratch: its proof requires
the full radical-extension theory. What we \emph{can} verify in Lean from first
principles is the group-theoretic half of the argument---that $S_5$ is not
solvable while $S_n$ ($n \le 4$) is---which is what the rest of this chapter does.
\end{warning}

Why? Adjoining an $n$th root corresponds to a cyclic (hence abelian) extension.
Solving by radicals means building a tower of radical extensions, each with a
cyclic Galois group. The overall Galois group must decompose into cyclic layers
---i.e., be solvable.

\section{$S_n$ for $n \leq 4$ Is Solvable}

\begin{prop}{}{s4-solvable}
$S_4$ is solvable, with composition series:
\[
  \{e\} \normal V_4 \normal A_4 \normal S_4
\]
where $V_4 = \{e, (12)(34), (13)(24), (14)(23)\}$ is the Klein four-group
and $A_4$ is the alternating group. The quotients are:
$V_4/\{e\} \iso V_4$ (abelian), $A_4/V_4 \iso \Z/3\Z$ (abelian),
$S_4/A_4 \iso \Z/2\Z$ (abelian).
\end{prop}

This is why the quartic formula exists: $|\Sym_4| = 24$, and $S_4$ has a
composition series with abelian quotients.

\section{$S_5$ Is Not Solvable}

\begin{thm}{$A_5$ Is Simple}{a5-simple}
The alternating group $A_5$ (order 60) has no nontrivial normal subgroups.
\end{thm}

\begin{proof}[Proof sketch]
$A_5$ has 60 elements partitioned into five conjugacy classes of sizes
$1, 12, 12, 15, 20$. A normal subgroup must be a union of conjugacy classes
including $\{e\}$ (size 1). The only sums of these sizes (including 1) that
divide 60 are 1 and 60 themselves. Therefore the only normal subgroups are
$\{e\}$ and $A_5$.
\end{proof}

\begin{cor}{}{s5-not-solvable}
$S_5$ is not solvable.
\end{cor}

\begin{proof}
If $S_5$ were solvable, then $A_5$ (as a subgroup) would also be solvable.
But $A_5$ is simple and non-abelian, so its only composition series is
$\{e\} \normal A_5$, with the non-abelian quotient $A_5/\{e\} = A_5$.
This contradicts solvability.
\end{proof}

\begin{thm}{The General Quintic Is Unsolvable}{quintic}
There is no formula in radicals for the roots of a general degree-5
polynomial.
\end{thm}

\begin{proof}
The generic degree-5 polynomial $f(x) = x^5 + a_4x^4 + \cdots + a_0$
over $\Q(a_0, \ldots, a_4)$ has Galois group $S_5$. Since $S_5$ is not
solvable, $f$ is not solvable by radicals.
\end{proof}

\begin{historical}[title={The Punchline}]
The jump from $S_4$ to $S_5$ is where simplicity (in the technical sense)
first appears. $A_4$ has the Klein four-group as a nontrivial normal subgroup;
$A_5$ has none. There is something fundamentally more complicated about
permuting 5 objects than permuting 4. Galois theory makes this precise:
the additional complexity is exactly what prevents a radical formula from
existing.
\end{historical}

\begin{exercisebox}[title={Exercise 10.1}]
\stepcounter{exercise}
Prove in Lean that $S_3$ is solvable by exhibiting the composition series
$\{e\} \normal \langle r \rangle \normal S_3$.
\end{exercisebox}

\begin{exercisebox}[title={Exercise 10.2}]
\stepcounter{exercise}
Prove that every group of order $\leq 4$ is solvable. \emph{Hint:} every
group of prime order is cyclic, hence abelian, hence solvable.
\end{exercisebox}

% ══════════════════════════════════════════════════════════════════════════════
\part{Applications and Connections}
% ══════════════════════════════════════════════════════════════════════════════

% ══════════════════════════════════════════════════════════════════════════════
\chapter{Finite Fields --- The Complete Classification}
% ══════════════════════════════════════════════════════════════════════════════

\epigraph{``Mathematics is not a careful march down a well-cleared highway, but
a journey into a strange wilderness.''}{\textit{W.S.\ Anglin}}

\section{Galois's Legacy}

\begin{historical}[title={The Fields That Bear His Name}]
In his final manuscripts, Galois described the finite fields: for every prime
power $p^n$, there is exactly one field of that size (up to isomorphism), and
these are \emph{all} the finite fields. They are denoted $\GF(p^n)$ (``Galois
field'') or $\F_{p^n}$ in his honor.
\end{historical}

\begin{thm}{Classification of Finite Fields}{finite-field-class}
\begin{enumerate}[nosep, label=(\roman*)]
\item For every prime $p$ and positive integer $n$, there exists a field of
  order $p^n$.
\item Any two fields of order $p^n$ are isomorphic.
\item Every finite field has order $p^n$ for some prime $p$ and some $n \geq 1$.
\end{enumerate}
\end{thm}

\section{Constructing Finite Fields}

\begin{defn}{}{gf-construction}
$\GF(p^n) = \F_p[x]/(f(x))$ where $f$ is any irreducible polynomial of degree
$n$ over $\F_p$.
\end{defn}

The key properties:
\begin{itemize}[nosep]
\item The multiplicative group $\GF(p^n)^\times$ is \textbf{cyclic} of order $p^n - 1$.
\item The \textbf{Frobenius automorphism} $\varphi : x \mapsto x^p$ generates
  $\Gal(\GF(p^n)/\F_p) \iso \Z/n\Z$.
\item The subfields of $\GF(p^n)$ correspond to divisors of $n$:
  $\GF(p^d)$ is a subfield iff $d \mid n$.
\end{itemize}

The subfield lattice is isomorphic to the divisor lattice of $n$---yet
another lattice appearing naturally.

\begin{exercisebox}[title={Exercise 11.1}]
\stepcounter{exercise}
Find an irreducible polynomial of degree 2 over $\F_3$ and use it to construct
$\GF(9)$. Write out the addition and multiplication tables.
\end{exercisebox}

% ══════════════════════════════════════════════════════════════════════════════
\chapter{$\GF(2^8)$ and the AES S-Box --- Capstone}
% ══════════════════════════════════════════════════════════════════════════════

\epigraph{``The enemy knows the system.''}{\textit{Claude Shannon}}

\section{AES and Galois Fields}

\begin{historical}[title={The Advanced Encryption Standard}]
The Advanced Encryption Standard (AES), adopted by NIST in 2001 after a
five-year international competition, encrypts the vast majority of the
world's internet traffic. Its designers---Joan Daemen and Vincent Rijmen
(Rijndael)---chose to build the cipher's core nonlinear operation around
arithmetic in $\GF(2^8)$ because:
\begin{itemize}[nosep]
\item Addition in $\GF(2^8)$ is XOR---free in hardware.
\item Multiplication can use lookup tables---fast in software.
\item The multiplicative inverse has excellent cryptographic properties.
\end{itemize}
\end{historical}

\section{Constructing $\GF(2^8)$}

We use the irreducible polynomial $m(x) = x^8 + x^4 + x^3 + x + 1$ over
$\F_2$:
\[
  \GF(2^8) = \F_2[x]/(x^8 + x^4 + x^3 + x + 1).
\]
Elements are bytes---polynomials of degree $\leq 7$ with coefficients in
$\{0, 1\}$. A byte like \texttt{0x53} represents $x^6 + x^4 + x + 1$.

\begin{leanbox}
\begin{minted}{lean}
-- GF(2⁸): elements are bytes, representing polynomials over F₂
structure GF256 where
  val : UInt8
  deriving DecidableEq, Repr

namespace GF256

-- Addition = XOR (polynomial addition mod 2)
def add (a b : GF256) : GF256 := ⟨a.val ^^^ b.val⟩

-- Multiplication = polynomial multiplication mod m(x)
-- m(x) = x⁸ + x⁴ + x³ + x + 1 = 0x11B
def mul (a b : GF256) : GF256 :=
  -- Russian peasant multiplication with reduction
  go a.val b.val 0
where
  go (a b acc : UInt8) : GF256 :=
    -- use `if h : b == 0` so the else-branch has a proof `b ≠ 0`,
    -- which the termination argument below needs
    if h : b == 0 then ⟨acc⟩
    else
      let acc' := if b &&& 1 != 0 then acc ^^^ a else acc
      let a'   := if a &&& 0x80 != 0
                  then (a <<< 1) ^^^ 0x1B  -- reduce mod m(x)
                  else a <<< 1
      go a' (b >>> 1) acc'
  termination_by b.toNat
  decreasing_by
    have hb : b ≠ 0 := by simpa using h
    have hbn : b.toNat ≠ 0 := fun hz => hb (by
      have := congrArg UInt8.ofNat hz; simpa using this)
    have hsr : (b >>> 1).toNat = b.toNat / 2 := by
      rw [UInt8.toNat_shiftRight, Nat.shiftRight_eq_div_pow]; rfl
    rw [hsr]; omega

-- Multiplicative inverse via extended Euclidean algorithm
-- (or: a⁻¹ = a^(254) since a^(255) = 1 in GF(2⁸)×)
def inv (a : GF256) : GF256 :=
  -- a^254 = a^(2+4+8+16+32+64+128)
  -- Compute by repeated squaring
  let a2   := mul a a
  let a4   := mul a2 a2
  let a8   := mul a4 a4
  let a16  := mul a8 a8
  let a32  := mul a16 a16
  let a64  := mul a32 a32
  let a128 := mul a64 a64
  mul (mul (mul (mul (mul (mul a2 a4) a8) a16) a32) a64) a128

-- The AES S-Box: inverse then affine transform
def sbox (a : GF256) : GF256 :=
  let b := if a.val == 0 then ⟨(0 : UInt8)⟩ else inv a
  affineTransform b
where
  affineTransform (b : GF256) : GF256 :=
    -- The affine transformation specified by AES
    -- b' = Ab + c where A is a specific 8×8 matrix over F₂
    -- and c = 0x63
    sorry -- Matrix multiplication over F₂

-- Sanity checks: the standard AES multiplication test vector, and that
-- `inv` really inverts. Both pass by kernel computation.
example : (mul ⟨0x57⟩ ⟨0x13⟩).val = 0xFE := by native_decide
example : (mul ⟨0x53⟩ (inv ⟨0x53⟩)).val = 1 := by native_decide

end GF256
\end{minted}
\end{leanbox}

\section{Proving Properties}

Properties we can verify:
\begin{itemize}[nosep]
\item The S-box is a bijection (256 distinct outputs).
\item No fixed points: $\text{sbox}(x) \neq x$ for all $x$.
\item No ``opposite fixed points'': $\text{sbox}(x) \neq x \oplus \texttt{0xFF}$.
\end{itemize}

These properties are what make AES resistant to linear and differential
cryptanalysis.

\begin{exercisebox}[title={Exercise 12.1}]
\stepcounter{exercise}
Prove that $\GF(2^8)$ satisfies the field axioms. The hardest part is
showing associativity of multiplication (polynomial multiplication mod $m(x)$).
For a finite type, \lc{native\_decide} can handle this.
\end{exercisebox}

\begin{exercisebox}[title={Exercise 12.2}]
\stepcounter{exercise}
Compute the full AES S-box table (all 256 entries) and verify it against the
published AES specification.
\end{exercisebox}

% ══════════════════════════════════════════════════════════════════════════════
\chapter{Compass and Straightedge --- Impossibility Proofs}
% ══════════════════════════════════════════════════════════════════════════════

\epigraph{``We must know. We will know.''}{\textit{David Hilbert}}

\section{Three Ancient Problems}

\begin{historical}[title={2000 Years of Frustration}]
Three construction problems haunted mathematicians from ancient Greece until
the 19th century:
\begin{enumerate}[nosep]
\item \textbf{Trisecting an angle:} Given an arbitrary angle, divide it into
  three equal parts.
\item \textbf{Doubling the cube:} Given a cube, construct another with
  exactly twice the volume.
\item \textbf{Squaring the circle:} Given a circle, construct a square with
  the same area.
\end{enumerate}
All three were to be done with compass and straightedge alone. Generations of
mathematicians attempted these constructions and failed. The Galois-theoretic
framework explains \emph{why}.
\end{historical}

\section{Constructible Numbers}

The key insight: every compass-and-straightedge construction corresponds to
solving equations of degree $\leq 2$. Intersecting two lines gives a linear
equation; intersecting a line and a circle, or two circles, gives a quadratic.
So each construction step extends the ``known'' field by degree at most 2.

\begin{thm}{Constructibility Criterion}{constructible}
A real number $\alpha$ is constructible if and only if there exists a tower
of field extensions
\[
  \Q = F_0 \subset F_1 \subset \cdots \subset F_n \ni \alpha
\]
with $[F_{i+1} : F_i] = 2$ for each $i$. In particular, $[\Q(\alpha):\Q]$
must be a power of $2$.
\end{thm}

\section{Trisecting $60°$}

Trisecting $60°$ requires constructing $\cos(20°)$. Using the triple angle
formula $\cos(3\theta) = 4\cos^3\theta - 3\cos\theta$ with $\theta = 20°$:
\[
  \cos(60°) = \frac{1}{2} = 4\cos^3(20°) - 3\cos(20°).
\]
So $\alpha = \cos(20°)$ is a root of $8t^3 - 6t - 1 = 0$, which is
irreducible over $\Q$ (by the rational root theorem: the only possible
rational roots are $\pm 1, \pm 1/2, \pm 1/4, \pm 1/8$, and none work).

Therefore $[\Q(\cos 20°):\Q] = 3$. Since 3 is not a power of 2, $\cos(20°)$
is not constructible, and a general angle cannot be trisected.

\section{Doubling the Cube}

Doubling a unit cube requires constructing $\sqrt[3]{2}$. Since $x^3 - 2$ is
irreducible over $\Q$, $[\Q(\sqrt[3]{2}):\Q] = 3$. Again, 3 is not a power
of 2: impossible.

\section{Squaring the Circle}

Requires constructing $\sqrt{\pi}$. But $\pi$ is \textbf{transcendental}
(not a root of any polynomial with rational coefficients), as proved by
Ferdinand von Lindemann in 1882. Since $\pi$ is not algebraic over $\Q$,
$[\Q(\pi):\Q]$ is infinite---certainly not a power of 2.

\begin{historical}[title={Gauss and the 17-gon}]
On the positive side: Gauss proved at age 19 (1796) that the regular 17-gon
is constructible. He was so proud of this that he reportedly asked for a
regular 17-gon to be inscribed on his tombstone. The criterion: a regular
$n$-gon is constructible iff $n = 2^k \cdot p_1 \cdots p_m$ where the $p_i$
are distinct Fermat primes ($p = 2^{2^j} + 1$). The known Fermat primes are
3, 5, 17, 257, and 65537.
\end{historical}

\begin{exercisebox}[title={Exercise 13.1}]
\stepcounter{exercise}
Prove in Lean that $8t^3 - 6t - 1$ has no rational roots by checking all
possibilities from the rational root theorem.
\end{exercisebox}

% ══════════════════════════════════════════════════════════════════════════════
\chapter{Error-Correcting Codes and Reed-Solomon}
% ══════════════════════════════════════════════════════════════════════════════

\section{Codes Over Finite Fields}

\begin{historical}[title={Reed and Solomon (1960)}]
Irving Reed and Gustave Solomon, at MIT Lincoln Laboratory, described a
class of error-correcting codes based on polynomial evaluation over finite
fields. Reed-Solomon codes protect data on CDs, DVDs, Blu-rays, QR codes,
deep-space probes (Voyager), RAID storage, and digital television. The
mathematics is pure Galois theory.
\end{historical}

The idea: encode a message of $k$ symbols as a polynomial $p(x)$ of degree
$< k$ over $\GF(q)$. Evaluate it at $n > k$ distinct points
$\alpha_1, \ldots, \alpha_n$. The codeword is
$(p(\alpha_1), \ldots, p(\alpha_n))$. Since a degree $< k$ polynomial is
determined by $k$ values, any $k$ of the $n$ evaluations suffice to recover
$p$---and hence the message. The redundancy ($n - k$ extra symbols) allows
correcting up to $\lfloor(n-k)/2\rfloor$ errors.

The algebraic key: the error-locator polynomial has roots in $\GF(q)$ that
identify the error positions. Finding these roots uses the structure of the
finite field---specifically, the fact that $\GF(q^n)$ contains all the
needed roots of unity.

\begin{exercisebox}[title={Exercise 14.1}]
\stepcounter{exercise}
Implement a simple Reed-Solomon encoder over $\GF(2^8)$: given a message of
$k$ bytes, produce $n$ evaluation points.
\end{exercisebox}

% ══════════════════════════════════════════════════════════════════════════════
\chapter{Galois Connections Revisited}
% ══════════════════════════════════════════════════════════════════════════════

\section{The Big Picture}

We introduced Galois connections abstractly in Chapter~3. Now we've seen the
original example: $\fixgrp$ and $\fixfld$ form an antitone Galois connection
between the lattice of intermediate fields and the lattice of subgroups.

\begin{thm}{}{gc-everywhere}
The following are all instances of Galois connections:
\begin{enumerate}[nosep, label=(\roman*)]
\item \textbf{Galois theory:} intermediate fields $\leftrightarrow$ subgroups.
\item \textbf{Algebraic geometry (Nullstellensatz):} affine varieties
  $\leftrightarrow$ radical ideals. $V(I)$ = common zeros of $I$;
  $I(V)$ = polynomials vanishing on $V$.
\item \textbf{Topology:} interior $\leftrightarrow$ closure. Open sets form a
  lattice; closed sets form a lattice; interior and closure connect them.
\item \textbf{Logic:} theories $\leftrightarrow$ models. $\text{Th}(M)$ =
  sentences true in models $M$; $\text{Mod}(T)$ = models satisfying theory $T$.
\item \textbf{Computer science:} abstract interpretation. Concrete semantics
  $\leftrightarrow$ abstract domains. The abstraction and concretization maps
  form a Galois connection.
\end{enumerate}
\end{thm}

\begin{intuition}
The pattern: whenever you have two ``worlds'' connected by maps that translate
between them, and the translations approximately invert each other (with a
controlled loss of information), you likely have a Galois connection. The
``closed'' elements on each side---those not affected by the round trip---form
the part of each world that is perfectly captured by the other.
\end{intuition}

\begin{exercisebox}[title={Exercise 15.1}]
\stepcounter{exercise}
Verify that the maps $V$ and $I$ from algebraic geometry satisfy the Galois
connection axioms (for a polynomial ring over an algebraically closed field).
\end{exercisebox}

% ══════════════════════════════════════════════════════════════════════════════
\chapter{Looking Back --- The Unity of Algebra}
% ══════════════════════════════════════════════════════════════════════════════

\section{The Story Arc}

We began with Lagrange's observation (1770) that the roots of a polynomial
have symmetries. We followed the thread through Gauss, Abel, and Galois to
the fundamental theorem: the solvability of an equation is controlled by its
symmetry group. We built every structure in Lean~4---groups, rings, fields,
polynomials, extensions, automorphisms---and proved the key theorems.

\section{What We Proved}

\begin{itemize}[nosep]
\item \textbf{The quintic is unsolvable}---not because we lack cleverness,
  but because $S_5$ is not solvable. The structure of the symmetry group
  makes a radical formula impossible.
\item \textbf{AES works}---because $\GF(2^8)$ has the algebraic properties
  (high nonlinearity of the inverse map) that resist cryptographic attacks.
\item \textbf{Angles can't be trisected}---because $\cos(20°)$ has
  degree 3 over $\Q$, and 3 is not a power of 2.
\item \textbf{Error-correcting codes work}---because finite field extensions
  provide the algebraic structure for polynomial interpolation and error
  location.
\end{itemize}

\section{The Galois-Theoretic Worldview}

The deepest lesson: \textbf{to understand an object, study its symmetries.}
This philosophy, implicit in Galois's work, was made explicit by Felix Klein
in his 1872 \textbf{Erlangen Program}: classify geometries by their symmetry
groups. It was taken to its ultimate form by Alexander Grothendieck in the
1960s, who reshaped algebraic geometry around symmetry and structure in ways
that are still being absorbed by mathematicians today.

\begin{historical}[title={Hilbert's 12th Problem and Beyond}]
In 1900, David Hilbert posed 23 problems for the 20th century. His 12th
problem asked for an explicit construction of all abelian extensions of
arbitrary number fields---a generalization of Galois theory that remains
open in full generality. The Langlands program, one of the deepest ongoing
research efforts in mathematics, can be seen as a vast generalization of
Galois theory connecting number theory, representation theory, and geometry.
Hilbert would have been amazed.
\end{historical}

\section{Connections Forward}

\begin{itemize}[nosep]
\item \textbf{Category theory:} groups, rings, and fields are all categories;
  homomorphisms are functors; Galois connections are adjunctions. Category
  theory is the language that unifies all of abstract algebra.
\item \textbf{Algebraic number theory:} Galois theory for number fields
  (extensions of $\Q$) is the foundation of modern number theory, including
  the proof of Fermat's Last Theorem (Wiles, 1995).
\item \textbf{Algebraic geometry:} Grothendieck's schemes generalize the
  Galois correspondence from fields to geometric objects, connecting algebra
  and geometry at the deepest level.
\item \textbf{Your other books:} The lattice/propagator book's Galois
  connections are the abstract skeleton of the fundamental theorem.
  The lambda calculus book's Curry-Howard correspondence is another instance
  of ``two worlds connected by structure-preserving maps.'' The concurrency
  book's bisimulation is a greatest fixed point---the same Knaster-Tarski
  theorem from lattice theory that underlies the closure operators here.
\end{itemize}

\vfill
\begin{center}
\Large\itshape\color{deepblue}
``The algebraic symmetries of the roots of an equation\\
contain the secret of its solvability.''\\[0.5cm]
\normalsize --- The message of Galois theory,\\
written in haste the night of May 29, 1832.
\end{center}
\vfill

% ══════════════════════════════════════════════════════════════════════════════
\appendix
% ══════════════════════════════════════════════════════════════════════════════

\chapter{A Lean 4 Primer}\label{app:primer}

This appendix is the promised zero-prerequisites tour of every Lean~4 construct
and tactic used in this book. Each entry says what the construct does, shows a
small self-contained example (every example here is machine-checked against
\texttt{leanprover/lean4:v4.31.0}, core Lean only), and---where it
illuminates---shows the lower-level term the tactic is sugar for. That last part
is worth your attention: Lean tactics are not magic, they are \emph{programs that
build proof terms}, and seeing the term once demystifies them permanently.
Everything here is core Lean; the one Mathlib-only tactic the book mentions
(\lc{ring}) is flagged as such.

\section{Declarations}

\subsection*{\lc{inductive} --- defining new types}

An \lc{inductive} declaration creates a brand-new type from a list of
\emph{constructors}---the only ways to build values of the type. The kernel
knows constructors are injective and disjoint, and it generates an induction
principle automatically (see \lc{induction} below). \lc{S3} in Chapter~1 and
\lc{GF256}'s wrapper are inductive/structure declarations.

\begin{leanbox}
\begin{minted}{lean}
inductive Tree where
  | leaf | node (l r : Tree)
  deriving Repr        -- auto-generate a printer, so #eval can show Trees
\end{minted}
\end{leanbox}

\lc{deriving Repr} asks Lean to synthesise a printer; \lc{deriving DecidableEq}
(used by \lc{S3}) synthesises a decidable \lc{==}.

\subsection*{\lc{def} and pattern matching}

A \lc{def} introduces a function or value. Functions on inductive types are
usually written by \emph{pattern matching}, one equation per constructor.
Recursion is allowed as long as some argument structurally shrinks; when Lean
cannot see why, you supply a \lc{termination\_by}/\lc{decreasing\_by} block (as
\lc{GF256.mul} does).

\begin{leanbox}
\begin{minted}{lean}
def Tree.size : Tree → Nat
  | .leaf     => 1
  | .node l r => l.size + r.size
\end{minted}
\end{leanbox}

The leading dot in \lc{.leaf} means ``the constructor of the type being
matched.'' Because the function is named \lc{Tree.size}, \emph{dot notation}
\lc{t.size} works for any \lc{t : Tree}.

\subsection*{\lc{structure}, \lc{class}, \lc{instance}, \lc{extends}}

A \lc{structure} bundles fields into a record; a \lc{class} is a structure meant
to be found by \emph{instance search} (an interface, like \lc{Group} or
\lc{Field}); an \lc{instance} registers an implementation. \lc{extends} makes
one class inherit another's fields (\lc{Field extends Ring}). The angle brackets
\lc{⟨…⟩} are the \emph{anonymous constructor}: Lean infers which constructor to
apply from the expected type. They also build proofs of $\exists$ and $\land$.

\begin{leanbox}
\begin{minted}{lean}
structure Point where
  x : Nat
  y : Nat
  deriving Repr

instance : Add Point where add a b := ⟨a.x + b.x, a.y + b.y⟩

example : (⟨1,2⟩ : Point) + ⟨3,4⟩ = ⟨4,6⟩ := rfl
\end{minted}
\end{leanbox}

This is exactly how \lc{Group} gets \lc{*} (Chapter~1) and \lc{Ring} gets \lc{+}
and \lc{*} (Chapter~4): we register \lc{Mul}/\lc{Add} instances that delegate to
the class fields.

\subsection*{\lc{theorem}, \lc{example}, and \lc{Prop}}

In Lean, propositions are types (of sort \lc{Prop}) and proofs are programs: to
prove $P$ is to build a term of type $P$. A \lc{theorem} is a named proof; an
\lc{example} is an anonymous one (used throughout this appendix). The body is
either a direct \emph{term} or a \emph{tactic block} introduced by \lc{by}.

\begin{leanbox}
\begin{minted}{lean}
example (p q : Prop) (hp : p) (h : p → q) : q := h hp   -- term style

example (p q : Prop) : p → (p → q) → q := by            -- tactic style
  intro hp h
  apply h
  exact hp
\end{minted}
\end{leanbox}

\section{Core Tactics}

\subsection*{\lc{rfl} --- true by computation}

Proves \lc{a = a} \emph{after both sides compute} to the same normal form.
Definitional unfolding is free, so \lc{2 + 3 = 5} is \lc{rfl}.

\begin{leanbox}
\begin{minted}{lean}
example : 2 + 3 = 5 := rfl
\end{minted}
\end{leanbox}

\subsection*{\lc{intro}, \lc{exact}, \lc{apply}}

If the goal is \lc{p → q} (or \lc{∀ x, P x}), \lc{intro hp} moves the antecedent
into the context. \lc{exact h} closes the goal with the term \lc{h}; \lc{apply h}
works \emph{backwards}: if \lc{h : p → q} and the goal is \lc{q}, it leaves goal
\lc{p}. The three tactics in the example above assemble exactly the term-style
proof \lc{h hp}.

\subsection*{\lc{rw} --- rewrite with an equation}

\lc{rw [h]}, for \lc{h : a = b}, replaces \lc{a} by \lc{b} in the goal
(\lc{rw [← h]} goes the other way, \lc{rw [h] at h2} rewrites a hypothesis).

\begin{leanbox}
\begin{minted}{lean}
example (a b : Nat) (h : a = b) : a + 1 = b + 1 := by rw [h]
\end{minted}
\end{leanbox}

\subsection*{\lc{simp} --- the simplifier}

\lc{simp} rewrites left-to-right with its rule set (lemmas tagged \lc{@[simp]},
plus anything you pass in \lc{simp [f, h]}) until nothing applies. Conceptually
it is \lc{rw} on autopilot; the \lc{conjugate\_*} lemmas of Chapter~7 are closed
by \lc{simp [conjugate]}.

\begin{leanbox}
\begin{minted}{lean}
example (n : Nat) : n + 0 = n := by simp
\end{minted}
\end{leanbox}

\subsection*{\lc{decide} and \lc{native\_decide}}

For a \emph{decidable} proposition, \lc{decide} runs the decision procedure in
the kernel and checks the answer is \lc{true}. \lc{native\_decide} does the same
but \emph{compiles} the check to native code---much faster on big goals (the
216-case $S_3$ associativity check, the AES vectors), at the cost of trusting the
compiler as well as the kernel. \lc{decide} can be slow on such goals.

\begin{leanbox}
\begin{minted}{lean}
example : (7 : Nat) < 10 := by decide
example : List.range 4 = [0,1,2,3] := by native_decide
\end{minted}
\end{leanbox}

\subsection*{\lc{omega} --- linear arithmetic}

\lc{omega} discharges goals in linear arithmetic over \lc{Nat} and \lc{Int}
(equalities and inequalities of $+$, $-$, constants, $<$, $\le$). It needs the
relevant facts to be in scope---the failure mode in Chapter~1's \lc{Fin.addMod}
is precisely \lc{omega} with no hypothesis proving \lc{0 < n}.

\begin{leanbox}
\begin{minted}{lean}
example (n : Nat) (h : n + 2 ≤ 5) : n ≤ 3 := by omega
\end{minted}
\end{leanbox}

\subsection*{\lc{cases} --- and the recursor underneath}

\lc{cases x} splits the goal into one branch per constructor of \lc{x}'s type.
It is sugar for applying the type's auto-generated \emph{recursor}. The two
proofs below are the same proof term:

\begin{leanbox}
\begin{minted}{lean}
example (b : Bool) : b = true ∨ b = false := by
  cases b with
  | true  => exact Or.inl rfl
  | false => exact Or.inr rfl

-- identical, written directly with the recursor Bool.rec:
example (b : Bool) : b = true ∨ b = false :=
  Bool.rec (Or.inr rfl) (Or.inl rfl) b
\end{minted}
\end{leanbox}

\subsection*{\lc{induction} --- recursion in proofs}

\lc{induction n} is \lc{cases} plus an \emph{induction hypothesis} \lc{ih} in
each recursive branch. It desugars to the recursor \lc{Nat.rec (motive := …)}
applied to the base case and the step (which takes \lc{ih}).

\begin{leanbox}
\begin{minted}{lean}
theorem zero_add (n : Nat) : 0 + n = n := by
  induction n with
  | zero => rfl
  | succ k ih => rw [Nat.add_succ, ih]   -- uses the hypothesis ih
\end{minted}
\end{leanbox}

\subsection*{\lc{obtain} --- destructuring, and \lc{Exists.elim}}

\lc{obtain ⟨w, hw⟩ := h} unpacks an existential or a conjunction, naming the
witness \lc{w} and its property \lc{hw}. It is sugar for \lc{Exists.elim}
(equivalently, a \lc{match}).

\begin{leanbox}
\begin{minted}{lean}
example (h : ∃ n : Nat, n = 3) : True := by
  obtain ⟨n, hn⟩ := h; trivial

example (h : ∃ n : Nat, n = 3) : True :=
  Exists.elim h (fun _ _ => trivial)
\end{minted}
\end{leanbox}

\subsection*{\lc{calc} --- chained (in)equalities, and \lc{Trans.trans}}

\lc{calc} writes a transitive chain step by step; each \lc{\_} is the previous
right-hand side. Under the hood the steps are glued with \lc{Trans.trans}.

\begin{leanbox}
\begin{minted}{lean}
example (a b c : Nat) (h1 : a = b) (h2 : b = c) : a = c := by
  calc a = b := h1
    _ = c := h2
\end{minted}
\end{leanbox}

\subsection*{Proof by contradiction --- \lc{Classical.byContradiction}}

\begin{sloppypar}
Mathlib provides a \lc{by\_contra} tactic. In core Lean the same move is the
term \lc{Classical.byContradiction}: it turns a proof of $\lnot\lnot p$ into a
proof of $p$ (classical double-negation elimination).
\end{sloppypar}

\begin{leanbox}
\begin{minted}{lean}
example (p : Prop) (h : ¬ ¬ p) : p :=
  Classical.byContradiction (fun hn => h hn)
\end{minted}
\end{leanbox}

\section{Termination and Other Constructs}

\subsection*{\lc{where}}

Attaches local auxiliary definitions to a \lc{def} (like \lc{go} inside
\lc{modPow} and \lc{GF256.mul}). They can see the outer parameters and are the
natural home for a recursive worker.

\subsection*{\lc{termination\_by} / \lc{decreasing\_by}}

When structural recursion is not obvious, \lc{termination\_by e} tells Lean which
expression \lc{e} decreases, and \lc{decreasing\_by tac} proves it strictly
decreases at each recursive call. \lc{GF256.mul} (Chapter~12) uses
\lc{termination\_by b.toNat} and a \lc{decreasing\_by} block rewriting
\lc{(b >>> 1).toNat} to \lc{b.toNat / 2} and closing with \lc{omega}.

\subsection*{\lc{ring} (Mathlib) vs.\ core Lean}

\lc{ring} proves any identity that holds in every commutative ring by
normalising both sides. It is \textbf{not} a core tactic---it lives in
Mathlib---so this book, which stays in core Lean, discharges such goals by hand
with the relevant \lc{Rat}/\lc{Int} lemmas (see \lc{conjugate\_mul} in
Chapter~7). If you \lc{import Mathlib}, a single \lc{ring} replaces that
bookkeeping.

\subsection*{\lc{sorry}}

A placeholder that proves \emph{anything}, with a warning. We use it only for
the handful of exercises left to the reader (e.g.\ \lc{GroupHom.map\_one}) and
for the two deep theorems (Chapters~8 and~10) we state but do not build from
scratch. A finished development has no \lc{sorry}.

\chapter{Lean 4 Typeclass Hierarchy}

The algebraic structures in this book form a hierarchy:

\begin{center}
\begin{tikzpicture}[
    node distance=1.5cm,
    box/.style={draw, rounded corners, fill=lightblue!60, font=\small\ttfamily,
      minimum width=2.5cm, minimum height=0.7cm},
    arr/.style={->, thick},
  ]
  \node[box] (po) {PartialOrder};
  \node[box, right=3cm of po] (grp) {Group};
  \node[box, above=1cm of po] (lat) {Lattice};
  \node[box, above=1cm of grp] (ring) {Ring};
  \node[box, above=1cm of ring] (field) {Field};
  \node[box, left=0.5cm of grp] (subgrp) {Subgroup};

  \draw[arr] (po) -- (lat);
  \draw[arr] (grp) -- (ring);
  \draw[arr] (ring) -- (field);
  \draw[arr] (grp) -- (subgrp);
  \draw[arr, dashed] (subgrp) -- node[above, font=\scriptsize\itshape] {forms a} (lat);
\end{tikzpicture}
\end{center}

\begin{itemize}[nosep]
\item \lc{Group} --- one operation with associativity, identity, inverses.
\item \lc{Ring} --- two operations; additive group + multiplicative monoid + distributivity.
\item \lc{Field} --- commutative ring where every nonzero element is invertible.
\item \lc{PartialOrder} --- reflexive, antisymmetric, transitive relation.
\item \lc{Lattice} --- partial order with meet and join.
\item \lc{Subgroup G} --- the subgroups of any group form a \lc{Lattice}.
\item \lc{GaloisConnection} --- order-reversing adjunction between two posets.
\end{itemize}

\end{document}
